
% BEGIN COMPLETE PRECISION TRANSFER AND PHASE PROOFS

\clearpage\begingroup
\part*{Actual arithmetic moment precision and all inverse margins}
\newtheorem{mprectheorem}{Theorem}
\newtheorem{mpreclemma}[mprectheorem]{Lemma}
\newtheorem{mprecproposition}[mprectheorem]{Proposition}
\newcommand{\mprecC}{\mathbb C}
\newcommand{\mprecR}{\mathbb R}
\newcommand{\mprectr}{\operatorname{tr}}
\newcommand{\mprecdiag}{\operatorname{diag}}
\newcommand{\mprecnorm}[1]{\left\lVert#1\right\rVert}


This calculation supplies quantitative conditioning and precision
constants for the actual finite-moment receiver in \cite{mprec:FM} and the
deflation-aware rank-one refinement RR1--16 in \cite{mprec:RR}. Both complete
proofs were read for this derivation. Positivity is proved from the
original native envelope, rather than assumed of a numerical midpoint.
All monomial coordinates, masses, parity blocks, complex transport,
observation and conductor fibres, and signed endpoint errors remain
unchanged. The elementary quadrature context is the human work of
Zimmerling, Druskin and Simoncini \cite{mprec:ZDS}; no finite-matrix
nondeflation hypothesis is imported from it.

\section{The original measure and explicit interval floor}

Retain the stipulated quartet domain and every original constant:
\begin{equation}\label{mprec:eq:objects}
\begin{gathered}
 \rho_*=\tfrac12+\delta+i\gamma,\quad 0<\delta<\tfrac12,
 \quad\gamma>2,\quad m\ge1,\quad g=2\xi,\\
 h(s)=\prod_{z\in\{\rho_*,\bar\rho_*,1-\rho_*,1-\bar\rho_*\}}(s-z)^m,
 \quad w_h(y)=\frac{|(g/h)(1/2+iy)|^2}{2\pi},\\
 m_k=w_h^{*k},\quad k=4\ell+1\ge9,\quad
 c=k/2,\quad S=c+iy,\quad \mu_h=\int_\mprecR w_h(y)dy,\\
 e_k=1+k(m-1),\quad q=e_k(k+1)^2,\quad
 \chi(S)=\prod_{a,b=0}^k
 [S-c-(2a-k)\delta-i(2b-k)\gamma]^{e_k}.
\end{gathered}
\end{equation}
This is the existing quartet domain, not an existence assertion for
an off-line zero. The complete original provider \cite{mprec:CAI},
CAI16--17, and its multiplicity continuation \cite{mprec:HG}, HG21, give
\begin{equation}\label{mprec:eq:envelope}
\begin{gathered}
 a_k(1+y^2)^{-M}\rho_\sigma(y)\le m_k(y)\le A_k\rho_\sigma(y),
 \quad \rho_\sigma(y)=\frac{|\Gamma(1/4+iy/2)|^2}{2\pi},\\
 c_\Gamma e^{-\alpha|y|}(1+|y|)^{-1/2}
 \le\rho_\sigma(y)\le
 C_\Gamma e^{-\alpha|y|}(1+|y|)^{-1/2},\\
 \alpha=\pi/2,\quad M=21+4m,\quad p=8m-9/2,\quad B_h=42+8m,\\
 c_\Gamma=\sqrt2e^{-7/6},\qquad C_\Gamma=\sqrt{2\sqrt5}e^{2/3},\\
 \vartheta_h=\int_{-1}^1w_h(y)dy,\qquad
 M_\alpha=\int_\mprecR e^{\alpha y}w_h(y)dy,\\
 a_k=\frac{c_h\vartheta_h^{k-3}e^{-\alpha(k-3)}(k-2)^{-B_h}}
 {C_\Gamma2^{B_h/2}},\qquad
 A_k=\frac{C_h^{\rm tilt}}{c_\Gamma}k^{p+1}M_\alpha^{k-1}.
\end{gathered}
\end{equation}
The constants $c_h,C_h^{\rm tilt}$ are exactly the provider's TW7
lower-envelope constant and BT12 compact/tail maximum. They and the
displayed original scalar integrals are not assigned numerical values
in this note. The Gamma density is an envelope only, never a
replacement measure.

The original parity and tilted measures are
\begin{equation}\label{mprec:eq:measures}
 d\nu_0(x)=m_k(\sqrt x)x^{-1/2}dx,\quad d\nu_1=x\,d\nu_0,
 \quad s(x)=\frac{\tanh(\alpha\sqrt x)}{\sqrt x},\quad
 d\widetilde\nu_\epsilon=s\,d\nu_\epsilon\quad(\epsilon=0,1).
\end{equation}
The original positive Stieltjes identity \cite{mprec:PP,mprec:FM} is
\begin{equation}\label{mprec:eq:stieltjes}
 s(x)=\frac4\pi\sum_{r=0}^\infty\frac1{x+(2r+1)^2},\qquad
 \frac1{s(x)}=\frac2\pi+\frac4\pi\sum_{j=1}^\infty\frac{x}{x+4j^2},
 \qquad 0<s(x)\le\alpha.
\end{equation}
Evenness and $x=y^2$ give, without any mass change,
\[
 \int|E(y^2)+yO(y^2)|^2m_k(y)dy
 =\int|E|^2d\nu_0+\int|O|^2d\nu_1.
\]
In particular $\nu_0((0,\infty))=\mu_h^k$ and
$\nu_1((0,\infty))=k\eta_2\mu_h^{k-1}$, where
$\eta_2=\int y^2w_h(y)dy$. The tilted masses remain their actual
integrals, and $\widetilde\nu_1=x\widetilde\nu_0$.

\begin{mpreclemma}[An actual native floor]\label{mprec:lem:floor}
Put
\begin{equation}\label{mprec:eq:densityfloor}
 d_\nu=\frac{a_kc_\Gamma3^{-M}e^{-\alpha\sqrt2}}
 {\sqrt2\sqrt{1+\sqrt2}},\qquad
 d_{\widetilde\nu}=\frac4{3\pi}d_\nu.
\end{equation}
For either parity, the density of $\nu_\epsilon$ on $[1,2]$ is at
least $d_\nu>0$, and that of $\widetilde\nu_\epsilon$ is at least
$d_{\widetilde\nu}>0$.
\end{mpreclemma}
\begin{proof}
On this interval $1+x\le3$, $\sqrt x\le\sqrt2$, and
$1+\sqrt x\le1+\sqrt2$. Substitution in \eqref{mprec:eq:envelope}
and the factor $x^{-1/2}$ gives the first floor for $\nu_0$;
the additional factor $x\ge1$ gives it for $\nu_1$.
The first positive term of \eqref{mprec:eq:stieltjes} gives
$s(x)\ge4/[\pi(x+1)]\ge4/(3\pi)$, which proves the second.
No interval-restricted measure replaces an original measure: we
only lower-bound its nonnegative quadratic integral by one part.
\end{proof}

For later upper margins define, for $\epsilon=0,1$ and $r\ge0$,
\begin{equation}\label{mprec:eq:momentcaps}
 U_{\nu_\epsilon,r}=
 \frac{2A_kC_\Gamma[2(r+\epsilon)]!}
 {\alpha^{2(r+\epsilon)+1}},\qquad
 U_{\widetilde\nu_\epsilon,r}=\alpha U_{\nu_\epsilon,r}.
\end{equation}
Then $\int x^r d\lambda\le U_{\lambda,r}$ for each of these
four actual measures. Indeed the substitution in
\eqref{mprec:eq:measures} gives $\int y^{2(r+\epsilon)}m_k(y)dy$;
discard $(1+|y|)^{-1/2}\le1$ in the upper envelope and integrate
$2\int_0^\infty y^Le^{-\alpha y}dy=2L!/\alpha^{L+1}$.
The tilted bound uses $s\le\alpha$.

\section{An exact rational Gram bound in the unchanged monomials}

For $t\ge0$, set $d=t+1$ and
\begin{equation}\label{mprec:eq:J}
 J_t=\left[\frac{2^{i+j+1}-1}{i+j+1}\right]_{i,j=0}^t,
 \quad D_t^{\rm int}=\det J_t,
 \quad T_t^{\rm int}=\mprectr J_t,
 \quad g_t=\frac{D_t^{\rm int}}{(T_t^{\rm int})^t}.
\end{equation}
All these are explicitly evaluated rational expressions in $t$.
The convention is $g_0=1$.

\begin{mpreclemma}\label{mprec:lem:J}
The matrix $J_t$ is positive definite and
\begin{equation}\label{mprec:eq:Jfloor}
 J_t\succeq g_tI,\qquad
 D_t^{\rm int}=\prod_{j=0}^t
 \frac{(j!)^4}{(2j)!(2j+1)!}
 =\prod_{j=0}^t\frac1{(2j+1)\binom{2j}j^2}.
\end{equation}
In particular $g_0=1$, $g_1=1/40$, and $g_2=5/981552$.
An entirely elementary lower estimate valid for all $t$ is
\begin{equation}\label{mprec:eq:gsize}
 g_t\ge(3/8)^t(2t+1)^{-(t+1)}4^{-(2t^2+t)}.
\end{equation}
\end{mpreclemma}
\begin{proof}
For a nonzero coefficient column $z$,
$z^*J_tz=\int_1^2|\sum_{j=0}^tz_jx^j|^2dx>0$;
a nonzero polynomial cannot vanish on an interval. If the eigenvalues
are $0<\lambda_1\le\cdots\le\lambda_d$, then
$\det J_t=\lambda_1\prod_{j=2}^d\lambda_j
\le\lambda_1(\mprectr J_t)^{d-1}$, proving the floor.

For the determinant identity alone put $x=1+u$. The coefficient
matrix $B_{lj}=\binom jl$ for $l\le j$ has determinant one, and
$J_t=B^*H_t^{\rm Hil}B$, where
$H_t^{\rm Hil}=[1/(i+j+1)]$. This auxiliary determinant identity
does not replace any native coefficient coordinates. The Cauchy
determinant identity yields
\[
 \det H_t^{\rm Hil}
 =\frac{\prod_{0\le i<j\le t}(j-i)^2}
 {\prod_{i,j=0}^t(i+j+1)}.
\]
For completeness, multiply the determinant of $[1/(u_i+v_j)]$
by $\prod_{i,j}(u_i+v_j)$. The resulting polynomial is alternating
in each list, hence divisible by the two Vandermonde products;
its total degree equals their combined degree. The quotient is a
constant, equal to one by the one-dimensional case and induction
on the residue at $u_t=-v_t$. Taking $u_i=i$, $v_j=j+1$ proves
the displayed formula. Grouping successive factors as $t$ grows
gives the product in \eqref{mprec:eq:Jfloor}. Direct substitution gives
the three small values. Finally $\binom{2j}j\le4^j$ implies
$D_t^{\rm int}\ge(2t+1)^{-(t+1)}4^{-t(t+1)}$, while
$T_t^{\rm int}\le2\sum_{i=0}^t4^i\le(8/3)4^t$.
Division proves \eqref{mprec:eq:gsize}.
\end{proof}

\section{Every required weighted native inverse has an explicit margin}

Fix one of the actual measures $\lambda$ in \eqref{mprec:eq:measures},
write $d_\lambda=d_\nu$ or $d_{\widetilde\nu}$, and put
$\mu_r=\int x^r d\lambda$. For a real polynomial
$w(x)=\sum_{l=0}^Lw_lx^l$ with $w_l\ge0$ and $w(1)>0$, set
\begin{equation}\label{mprec:eq:W}
 W_t(w)=\left[\sum_{l=0}^Lw_l\mu_{i+j+l}\right]_{i,j=0}^t,
 \quad h_{\lambda,t}=d_\lambda g_t,
 \quad \ell_t(w)=w(1)h_{\lambda,t},\quad
 U_t(w)=\sum_{i=0}^t\sum_{l=0}^Lw_lU_{\lambda,2i+l}.
\end{equation}
On $[1,2]$, $w(x)\ge w(1)$, so the proved native density floor
and \eqref{mprec:eq:Jfloor} give the full quantitative bounds
\begin{equation}\label{mprec:eq:Wfloor}
 \boxed{\ell_t(w)I\preceq W_t(w)\preceq U_t(w)I.}
\end{equation}
The upper bound follows from positivity and the trace cap. In
particular the moment systems actually used in \cite{mprec:FM} have
the following constants, not merely existential eigenvalues:
\[
\begin{array}{c|c|c}
 w(x)&\text{matrix}&\ell_t(w)/h_{\lambda,t}\ \ (=\sum_lw_l)\\ \hline
 x^l&H_t^{(l)}=[\mu_{i+j+l}]&1\\
 x+a&H_t^{(1)}+aH_t&1+a\\
 x(x+a)&H_t^{(2)}+aH_t^{(1)}&1+a\\
 x(x+a)(x+b)&H_t^{(3)}+(a+b)H_t^{(2)}+abH_t^{(1)}&(1+a)(1+b)
\end{array}
\]
Here $a,b>0$; in the original calculation $a=4j^2$ and
$b=(2r+1)^2$. These formulas apply separately to both parity
blocks, both with and without the actual factor $s$.

\begin{mprectheorem}[A pole-independent moment budget]\label{mprec:thm:budget}
Suppose each needed scalar moment has a real certified midpoint
$\widehat\mu_r$ and absolute error at most $\varepsilon$.
Let $\widehat W$ be its moment matrix. For $d=t+1$,
\begin{equation}\label{mprec:eq:deltaW}
 \mprecnorm{W-\widehat W}\le\delta_W:=d\,w(1)\varepsilon.
\end{equation}
For every $0<\theta\le1/4$, the concrete common budget
\begin{equation}\label{mprec:eq:uniformbudget}
 \boxed{\varepsilon\le\frac{\theta d_\lambda g_t}{t+1}}
\end{equation}
ensures $\delta_W\le\theta\ell_t(w)$ for every weight in
the table, simultaneously for all positive poles. In particular
\begin{equation}\label{mprec:eq:inversebracket}
 \widehat W-\delta_W I\succeq(\ell_t(w)-2\delta_W)I>0,
 \quad
 (\widehat W+\delta_W I)^{-1}\preceq W^{-1}
 \preceq(\widehat W-\delta_W I)^{-1},
\end{equation}
and
\begin{equation}\label{mprec:eq:inverseerror}
 \mprecnorm{W^{-1}-\widehat W^{-1}}
 \le\frac{\delta_W}{\ell_t(w)(\ell_t(w)-\delta_W)}.
\end{equation}
\end{mprectheorem}
\begin{proof}
Each entry error is at most $w(1)\varepsilon$. Its maximum
absolute row sum and hence the norm of the Hermitian error are
at most $d w(1)\varepsilon$. Weyl's elementary quadratic-form
inequality gives $\widehat W\succeq(\ell_t(w)-\delta_W)I$.
This proves the first assertion in \eqref{mprec:eq:inversebracket};
the other assertions use inversion of positive definite Loewner
inequalities and
$W^{-1}-\widehat W^{-1}=W^{-1}(\widehat W-W)\widehat W^{-1}$.
The coefficient sum and the interval minimum are the same $w(1)$,
so it cancels in \eqref{mprec:eq:uniformbudget}.
\end{proof}

For different moment radii $\varepsilon_r$, a sharper completely
explicit choice is
\[
 \delta_W=\max_{0\le i\le t}\sum_{j=0}^t
                \sum_lw_l\varepsilon_{i+j+l}.
\]
For a finite list of degrees and measures, take the minimum of the
right side of \eqref{mprec:eq:uniformbudget}. This is a specified finite
positive number, not a request to test an unproved midpoint
positivity. When a common base sequence represents both parities,
use the literal shift $\mu_{1,r}=\mu_{0,r+1}$, preserving its
correlation and mass.

\section{Outward quadratic expressions and explicit absolute budgets}

Every inverse-derived matrix in \cite{mprec:FM} has the form
$F=C^*W^{-1}C$, where $C=[\mu_{i+j+l_0}]$ is $d$ by $e=n+1$,
$l_0=0$ or $1$, and $d=t+1\ge e$. Define explicit numbers
\begin{equation}\label{mprec:eq:P}
 \begin{gathered}
 c_C=\sqrt{de},\qquad
 U_C=\left(\sum_{i=0}^{d-1}\sum_{j=0}^{e-1}
                     U_{\lambda,i+j+l_0}^{2}\right)^{1/2},\\
 c_W=d\,w(1),\quad \ell=\ell_t(w),\qquad
 P_F=\frac43\left\{\frac{c_WU_C^2}{\ell^2}
                  +\frac{2U_Cc_C+c_C^2}{\ell}\right\}.
 \end{gathered}
\end{equation}
These are evaluated algebraic expressions in the original upper
caps and lower floors. For $\varepsilon\le1$ and
$c_W\varepsilon\le\ell/4$, set
$\widehat F=\widehat C^*\widehat W^{-1}\widehat C$.
Then
\begin{equation}\label{mprec:eq:Ferror}
 \mprecnorm{F-\widehat F}\le P_F\varepsilon,
 \qquad \widehat F-P_F\varepsilon I\preceq F
 \preceq\widehat F+P_F\varepsilon I.
\end{equation}
Indeed $\mprecnorm{C}\le U_C$, $\mprecnorm{C-\widehat C}\le c_C\varepsilon$,
$\mprecnorm{\widehat W^{-1}}\le4/(3\ell)$ and
$\mprecnorm{W^{-1}-\widehat W^{-1}}\le4c_W\varepsilon/(3\ell^2)$.
Insert these in
\[
 F-\widehat F=C^*(W^{-1}-\widehat W^{-1})C
              +C^*\widehat W^{-1}C-
                    \widehat C^*\widehat W^{-1}\widehat C.
\]
The second difference has norm at most
$(2U_Cc_C\varepsilon+c_C^2\varepsilon^2)4/(3\ell)$;
$\varepsilon^2\le\varepsilon$ proves the claim. Thus a prescribed
absolute matrix error $\tau>0$ is attained with the explicit budget
\begin{equation}\label{mprec:eq:budgetF}
 \varepsilon\le\min\{1,\ell/(4c_W),\tau/P_F\}.
\end{equation}
The restriction $\varepsilon\le1$ selects an absolute accuracy;
it is not a normalization of a measure or of its moments.

There is also a native lower floor for the exact $F$ itself. The
first $e$ rows of $C$ form $H_n^{(l_0)}\succeq h_{\lambda,n}I$,
and $W^{-1}\succeq U_t(w)^{-1}I$. Consequently
\begin{equation}\label{mprec:eq:Ffloor}
 F\succeq\frac{h_{\lambda,n}^2}{U_t(w)}I.
\end{equation}
Choosing $P_F\varepsilon\le h_{\lambda,n}^2/[4U_t(w)]$
makes the lower outward endpoint in \eqref{mprec:eq:Ferror} strictly
positive with at least half this floor. This includes the
quadratic-denominator matrix with $l_0=1$ and
$w=x(x+a)(x+b)$, not only the one-pole systems.

An actual floating computation of the midpoint inverse must itself
be certified. One option is exact rational inversion, since the
moment midpoints may be rational and the original poles are positive
integers. More generally, if a computed matrix $R_0$ satisfies a
certified $\mprecnorm{I-\widehat W R_0}\le r_0<1$, then
\begin{equation}\label{mprec:eq:residualinverse}
 \mprecnorm{\widehat W^{-1}-R_0}
 \le\frac{\mprecnorm{R_0}r_0}{1-r_0}.
\end{equation}
To prove this put $E=I-\widehat W R_0$ and use
$\widehat W^{-1}=R_0(I-E)^{-1}$ and its convergent geometric
series. Replacing $R_0$ by its Hermitian part does not increase
the error from the exact Hermitian inverse. Its extra quadratic
error is at most $\mprecnorm{\widehat C}^2$ times the right side of
\eqref{mprec:eq:residualinverse}, plus certified multiplication roundoff.
All subsequent scalar coefficients, square roots and logarithms are
likewise rounded outwards if evaluated numerically.

\section{A quantitative Radau margin and the exact complex phase}

For the actual tilted measure $\rho=\widetilde\nu_\epsilon$,
use the original RR notation, with $r\ge n+1$:
\begin{equation}\label{mprec:eq:RR}
 \begin{gathered}
 H=[\mu_{i+j}]_{0\le i,j<r},\quad
 D=[\mu_{i+j+1}]_{0\le i,j<r},\quad
 C=[\mu_{i+j}]_{\substack{0\le i<r\\0\le j\le n}},\quad
 e=(0,\ldots,0,1)^T,\\
 \kappa=(e^*D^{-1}e)^{-1},\quad \widetilde D=D-\kappa ee^*,
 \quad W_G=D+aH,\quad W_R=\widetilde D+aH.
 \end{gathered}
\end{equation}
Here $\widetilde D$ is semidefinite, not positive definite, and
is never inverted by itself. The attained constrained minimum
\begin{equation}\label{mprec:eq:kmin}
 \kappa=\min_{e^*z=1}z^*Dz
\end{equation}
has minimizer $\kappa D^{-1}e$; subtracting it from an arbitrary
point in the fibre proves $\widetilde D\succeq0$ and its exact
one-dimensional kernel. Set
\begin{equation}\label{mprec:eq:RRconstants}
 \begin{gathered}
 h=d_{\widetilde\nu}g_{r-1},\quad
 U_H=\sum_{i=0}^{r-1}U_{\rho,2i},\quad
 U_D=\sum_{i=0}^{r-1}U_{\rho,2i+1},\quad U_\kappa=U_{\rho,2r-1}.
 \end{gathered}
\end{equation}
We have
\begin{equation}\label{mprec:eq:radaufloor}
 H,D\succeq hI,\quad 0<\kappa\le D_{r-1,r-1}\le U_\kappa,
 \quad W_G\succeq(1+a)hI,\quad W_R\succeq ahI.
\end{equation}
Thus the apparent Schur denominator has the explicit native floor
\begin{equation}\label{mprec:eq:denominator}
 \boxed{1-\kappa e^*W_G^{-1}e\ge
                    \frac{ah}{U_D+aU_H}>0.}
\end{equation}
To check it, conjugate $W_R\succeq ahI$ by $W_G^{-1/2}$.
The resulting matrix is $I-\kappa uu^*$, $u=W_G^{-1/2}e$,
and is at least $ah(U_D+aU_H)^{-1}I$. Its least eigenvalue
is precisely the left side of \eqref{mprec:eq:denominator}.

This lower margin has a concrete moment tolerance. For moment
error $\varepsilon$ with $r\varepsilon<h$, positivity of
$\widehat D$ follows already from the interval floor, so define
the exact midpoint scalar $\widehat\kappa=(e^*\widehat D^{-1}e)^{-1}$.
The error bound $\mprecnorm{D-\widehat D}\le r\varepsilon$ gives
\[
 (1-r\varepsilon/h)D\preceq\widehat D
 \preceq(1+r\varepsilon/h)D.
\]
Minimizing on the same fibre in \eqref{mprec:eq:kmin} proves
\begin{equation}\label{mprec:eq:kerror}
 |\widehat\kappa-\kappa|\le
                  \frac{rU_\kappa}{h}\varepsilon.
\end{equation}
Therefore, for
$\widehat W_R=\widehat D+a\widehat H-\widehat\kappa ee^*$,
\begin{equation}\label{mprec:eq:radauprecision}
 \mprecnorm{W_R-\widehat W_R}\le c_R\varepsilon,
 \quad c_R=r(1+a)+\frac{rU_\kappa}{h},\quad
 \varepsilon\le\min\{1,h/(4r),ah/(4c_R)\}
\end{equation}
guarantees the same positive outward inverse brackets as
\eqref{mprec:eq:inversebracket}, with $\ell=ah$, $c_W=c_R$.
Formula \eqref{mprec:eq:P} with those constants certifies
$F^+=C^*W_R^{-1}C$. Its Gauss counterpart
$F^-=C^*W_G^{-1}C$ uses $\ell=(1+a)h$, $c_W=r(1+a)$.
This proves a finite precision budget even for the deflated Radau
system; no generic numerical positivity assumption was used.

For clarity, the rank-one phase assertion used in this receiver
also survives the outward calculation. Let $p_r$ be the actual
monic polynomial orthogonal for $\rho$ and let $z=v_n(-a)$.
Orthogonality of the Gauss residual gives, for every
$f(x)=v_n(x)^Tu$,
\[
 f-(x+a)q=-q_{r-1}p_r,\qquad q_{r-1}=-f(-a)/p_r(-a),
 \quad[q]=W_G^{-1}Cu.
\]
The value $p_r(-a)$ cannot vanish, since division by $x+a$
would contradict positivity of $\int(x+a)|p_r/(x+a)|^2d\rho$.
Expanding the variational square yields
\begin{equation}\label{mprec:eq:phase}
 F_n(a)=\int\frac{v_nv_n^*}{x+a}d\rho
       =F^-+\lambda zz^*,\qquad
 \lambda=\frac1{p_r(-a)^2}\int\frac{p_r(x)^2}{x+a}d\rho\ge0.
\end{equation}
The rank-one inverse identity and
$e^*W_G^{-1}C=-z^*/p_r(-a)$ give
\[
 F^+-F^-=\gamma zz^*,\qquad
 \gamma=\frac{\kappa}{p_r(-a)^2(1-\kappa e^*W_G^{-1}e)}.
\]
For the upper sign, let $t_{r-1}$ be monic orthogonal for $x\rho$.
Then $\int x|t_{r-1}|^2d\rho=\kappa$, and the Radau equations
give $f-(x+a)\widetilde q=-\widetilde q_{r-1}xt_{r-1}$ with
$\widetilde q_{r-1}=f(-a)/[a t_{r-1}(-a)]$.
Divide $\bar f(x)-\bar f(-a)$, then
$t_{r-1}(x)-t_{r-1}(-a)$, by $x+a$ and use orthogonality
against $x\rho$. The complete difference is
\[
 u^*(F^+-F_n)u=
 \frac{|f(-a)|^2}{a t_{r-1}(-a)^2}
                   \int\frac{x t_{r-1}(x)^2}{x+a}d\rho\ge0.
\]
The divided differences vanish for constants, so $r=1$ is included.
Thus $0\le\lambda\le\gamma$ with its original mass.

If \eqref{mprec:eq:Ferror} certifies midpoint errors $\eta_\pm$,
one may evaluate $\gamma$ stably without polynomial division:
\begin{equation}\label{mprec:eq:gammaout}
 \widehat\gamma=(\widehat F^+-\widehat F^-)_{00},\qquad
 \max\{0,\widehat\gamma-\eta_+-\eta_-\}\le\gamma
 \le\widehat\gamma+\eta_++\eta_-=: \gamma^{\rm up}.
\end{equation}
The exact matrix description is
\begin{equation}\label{mprec:eq:phaseout}
 F_n=\widehat F^-+E^-+\lambda zz^*,\qquad
 \mprecnorm{E^-}\le\eta_-,\quad E^-=(E^-)^*,\quad
 0\le\lambda\le\gamma^{\rm up}.
\end{equation}
In particular its complex cross error is
$u^*E^-v+\lambda\overline{f_u(-a)}f_v(-a)$, not an
independently assigned phase. The same $E^-$ and scalar $\lambda$
serve every entry. This distinction is retained by every congruence.

\section{Finite tilted and resolvent enclosures from these budgets}

Here are the precise outward substitutions in the previous native
formulas. They also specify the cost of every finite moment error.
For the chosen $\lambda$, target $n$, and auxiliary $t\ge n$,
put $H_n=[\mu_{i+j}]$, $C=[\mu_{i+j}]_{0\le i\le t,0\le j\le n}$,
$D_0=[\mu_{i+j+1}]_{0\le i\le t,0\le j\le n}$ and
\[
 L=C^*(H_t^{(1)}+aH_t)^{-1}C,\qquad
 V=D_0^*(H_t^{(2)}+aH_t^{(1)})^{-1}D_0.
\]
Complete squares with weights $x+a$ and $x(x+a)$ respectively
give
\begin{equation}\label{mprec:eq:variational}
 L\preceq R_n(a)\preceq(H_n-V)/a,\quad
 V\preceq B_n(a)\preceq H_n-aL,
\end{equation}
where $R_n=\int vv^*/(x+a)d\lambda$,
$B_n=\int xvv^*/(x+a)d\lambda$. Explicitly the two lower
errors on a vector $u$ are the minima over degree-$t$ polynomials
$q$ of $\int(x+a)|f/(x+a)-q|^2d\lambda$ and
$\int x(x+a)|f/(x+a)-q|^2d\lambda$, respectively.
Subtracting each lower bound from $H_n=aR_n+B_n$ proves the
upper bounds. Thus no sign in \eqref{mprec:eq:variational} is assumed.

Let $\eta_H=(n+1)\varepsilon$ and $\eta_L=P_L\varepsilon$,
$\eta_V=P_V\varepsilon$. Certified finite endpoints are
\begin{equation}\label{mprec:eq:outvar}
 \begin{aligned}
 R^{\rm lo}&=\widehat L-\eta_LI,&
 R^{\rm up}&=(\widehat H_n-\widehat V+(\eta_H+\eta_V)I)/a,\\
 B^{\rm lo}&=\widehat V-\eta_VI,&
 B^{\rm up}&=\widehat H_n-a\widehat L+(\eta_H+a\eta_L)I.
 \end{aligned}
\end{equation}
For Radau replace these by
$B^{\rm lo}=\widehat H_n-a\widehat F^+-(\eta_H+a\eta_+)I$
and
$B^{\rm up}=\widehat H_n-a\widehat F^-+(\eta_H+a\eta_-)I$.
They enclose the exact $H_n-aF^+$ and $H_n-aF^-$ endpoints.

When only untilted moments are used, keep the exact mixed-pole
integral against $\nu_\epsilon$:
\[
 T_n(a,b)=\int\frac{xvv^*}{(x+a)(x+b)}d\nu_\epsilon,
 \quad T^- =D_0^*W_t(x(x+a)(x+b))^{-1}D_0.
\]
Completing the square in the weight $x(x+a)(x+b)$ for the
function $f/[(x+a)(x+b)]$ proves $T^-\preceq T_n$.
Its lower outward matrix is $\widehat T^--P_T\varepsilon I$.
The exact partial fraction identity gives the upper matrix
\begin{equation}\label{mprec:eq:partial}
 T^{\rm up}=
 \begin{cases}
 [bR^{\rm up}(b)-aR^{\rm lo}(a)]/(b-a),&b>a,\\
 [aR^{\rm up}(a)-bR^{\rm lo}(b)]/(a-b),&a>b.
 \end{cases}
\end{equation}
Indeed $T_n=[bR_n(b)-aR_n(a)]/(b-a)$; the indicated upper
and lower signs follow from the sign of each coefficient. The
original even-square and odd-square poles are never equal.

For explicit scalar precision accounting, the centers for the
exact finite $R^-$ and $R^+$ expressions have error coefficients
$\eta_{R^-}(a)=\eta_L(a)$ and
$\eta_{R^+}(a)=(\eta_H+\eta_V(a))/a$. Consequently the
center for the finite upper expression in \eqref{mprec:eq:partial}
has radius
\[
 \eta_{T^+}(a,b)=
 \begin{cases}
 [b\eta_{R^+}(b)+a\eta_{R^-}(a)]/(b-a),&b>a,\\
 [a\eta_{R^+}(a)+b\eta_{R^-}(b)]/(a-b),&a>b,
 \end{cases}
 \qquad \eta_{T^-}=P_T\varepsilon.
\]
These are the sums of absolute signed coefficients, not
unsupported subtraction of error bounds.

For $b_r=(2r+1)^2$ and any integer $R\ge1$ put
\[
 \theta_R=\alpha-\frac4\pi\sum_{r=0}^{R-1}\frac1{b_r}>0,
 \qquad\theta_R\le\frac2{\pi(2R-1)}.
\]
The first equality uses the value of the positive series at zero;
the bound compares the decreasing omitted terms with the integral
of $(2u+1)^{-2}$ from $R-1$ to infinity. If
$H^{\rm up}=\widehat H_n+\eta_HI$, then the actual tilted
matrices are enclosed by
\begin{equation}\label{mprec:eq:tiltedout}
 \begin{aligned}
 \frac4\pi\sum_{r<R}R^{\rm lo}(b_r)
 &\preceq \int vv^*s\,d\nu_\epsilon
 \preceq\frac4\pi\sum_{r<R}R^{\rm up}(b_r)+\theta_RH^{\rm up},\\
 \frac4\pi\sum_{r<R}T^{\rm lo}(a,b_r)
 &\preceq \int\frac{xvv^*s}{x+a}d\nu_\epsilon
 \preceq\frac4\pi\sum_{r<R}T^{\rm up}(a,b_r)+\theta_RH^{\rm up}.
 \end{aligned}
\end{equation}
The omitted positive scalar kernels are at most $\theta_R$,
since $x/(x+a)\le1$. Integration proves every remainder in
\eqref{mprec:eq:tiltedout}. These are finite original-moment enclosures
with all scalar tails and signed coefficients retained.

The four finite centers on the two lines of \eqref{mprec:eq:tiltedout}
therefore have explicit radii
\[
 \begin{array}{ll}
 \eta_{K^-}=\dfrac4\pi\sum_{r<R}\eta_{R^-}(b_r),&
 \eta_{K^+}=\dfrac4\pi\sum_{r<R}\eta_{R^+}(b_r)+\theta_R\eta_H,\\[3pt]
 \eta_{\widetilde B^-}=\dfrac4\pi\sum_{r<R}\eta_{T^-}(a,b_r),&
 \eta_{\widetilde B^+}=\dfrac4\pi\sum_{r<R}\eta_{T^+}(a,b_r)
                                      +\theta_R\eta_H.
 \end{array}
\]
Every coefficient here is the displayed finite expression in
original constants, poles and degrees; no unknown inverse norm
is left inside a precision request.

\section{Both parity blocks, outward source bounds, and direct precision}

At an original degree $N\ge q-1$ put
$n_0=\lfloor N/2\rfloor$, $n_1=\lfloor(N-1)/2\rfloor$, and
\[
 H_N=\mprecdiag\left([\int x^{i+j}d\nu_0]_{0\le i,j\le n_0},
                [\int x^{i+j}d\nu_1]_{0\le i,j\le n_1}\right).
\]
Let $K$ and $K_x$ be the corresponding actual tilted moment
blocks with no shift and one shift. The original reciprocal series
in \eqref{mprec:eq:stieltjes}, the one-pole bounds above, and
$\sum_{j>J}j^{-2}\le1/J$ give
\begin{equation}\label{mprec:eq:Gexact}
 G^-:=\frac2\pi K+\frac4\pi\sum_{j=1}^JB_j^-
 \preceq H_N\preceq
 G^+:=\frac2\pi K+\frac4\pi\sum_{j=1}^JB_j^+
                        +\frac1{\pi J}K_x.
\end{equation}
Here $B_j^\pm$ may be the exact finite variational or Radau
endpoints in both blocks. They are positive semidefinite: for
Radau, $W_R\succeq aH$ gives
$aC^*W_R^{-1}C\preceq C^*H^{-1}C=H_n$; the variational
lower endpoint is manifestly positive. Therefore
\begin{equation}\label{mprec:eq:gsource}
 G^-\succeq g_N^*I,\qquad
 g_N^*=\frac2\pi d_{\widetilde\nu}
                         \min\{g_{n_0},g_{n_1}\}>0.
\end{equation}
Apply the moment budgets to every displayed finite matrix. If
its exact endpoint has a center $\widehat G^\pm$ with certified
norm error $\eta_\pm$, set
\begin{equation}\label{mprec:eq:Gout}
 G^{\rm lo}=\widehat G^--\eta_-I,\qquad
 G^{\rm up}=\widehat G^++\eta_+I.
\end{equation}
For the direct tilted route one may take, explicitly,
\[
 \eta_-\le\frac2\pi\eta_K+\frac4\pi\sum_{j=1}^J\eta_{B_j^-},
 \quad
 \eta_+\le\frac2\pi\eta_K+\frac4\pi\sum_{j=1}^J\eta_{B_j^+}
                         +\frac1{\pi J}\eta_{K_x}.
\]
Each $\eta$ on the right is the displayed coefficient times
$\varepsilon$, or the maximum of the two parity coefficients.
Thus these give finite known coefficients $P_\pm$ with
$\eta_\pm\le P_\pm\varepsilon$; choose
$\varepsilon\le g_N^*/(4P_-)$ as well as all local budgets.
Then $G^{\rm lo}\succeq(g_N^*-2\eta_-)I\succeq g_N^*I/2$,
and
\begin{equation}\label{mprec:eq:sourcecert}
 0<G^{\rm lo}\preceq H_N\preceq G^{\rm up},\quad
 \Delta G=G^{\rm up}-G^{\rm lo}\succeq0,\quad
 G^{\rm up}\preceq(1+\eta_N)G^{\rm lo},\quad
 \eta_N=\frac{\mprecnorm{\Delta G}}{g_N^*-2\eta_-}.
\end{equation}
The exact norm may be replaced by a certified Frobenius or row-sum
upper bound. The original-moment-only route \eqref{mprec:eq:tiltedout}
uses the same signed scalar error addition, with the actual native
outer tail $X_H/(2J)$ in place of $K_x/(\pi J)$, where $X_H$
is the once-shifted untilted direct sum. This follows from
$s\le\pi/2$, with no substituted measure. In that route the
exact finite lower approximation to $K$ is smaller than $K$,
so its positivity floor must also be replaced. Keep its first
pole $b_0=1$ and let $t_{\epsilon,0}$ be that pole's actual
auxiliary degree. Formula \eqref{mprec:eq:Ffloor} gives the explicit
valid replacement
\[
 g_{N,\mathrm{untilted}}^*
 =\frac8{\pi^2}\min_{\epsilon=0,1}
 \frac{(d_\nu g_{n_\epsilon})^2}
 {U_{t_{\epsilon,0}}^{(\nu_\epsilon)}(x+1)}>0.
\]
All other finite lower summands are positive. Use this floor,
the literal signed coefficients from \eqref{mprec:eq:partial}, and
the tails in \eqref{mprec:eq:tiltedout} in
\eqref{mprec:eq:Gout}--\eqref{mprec:eq:sourcecert}; the same proved
$\eta_-\le g_{N,\mathrm{untilted}}^*/4$ budget certifies the
lower source endpoint. No positivity floor for the exact tilted
Gram is mistakenly assigned to its finite lower approximation.

There is a stronger elementary route when the original moments
through degree $2n_\epsilon$ are supplied: they already give
$H_N$ directly. There is no need to reconstruct it by resolvents.
Define
\begin{equation}\label{mprec:eq:directconstants}
 h_N=d_\nu\min\{g_{n_0},g_{n_1}\},\qquad
 d_N^{\rm src}=\max\{n_0+1,n_1+1\},\qquad
 \delta_N=d_N^{\rm src}\varepsilon.
\end{equation}
Then $H_N\succeq h_NI$ and $\mprecnorm{H_N-\widehat H_N}\le\delta_N$.
For $\delta_N\le\theta_Nh_N$, $0<\theta_N\le1/4$, put
\begin{equation}\label{mprec:eq:directout}
 H_N^{\rm lo}=\widehat H_N-\delta_NI,\quad
 H_N^{\rm up}=\widehat H_N+\delta_NI.
\end{equation}
They are positive outward endpoints and obey
\begin{equation}\label{mprec:eq:directrelative}
 H_N^{\rm up}\preceq\frac1{1-2\theta_N}H_N^{\rm lo}.
\end{equation}
Indeed $H_N^{\rm lo}\succeq(h_N-2\delta_N)I$ and the two
endpoints differ by $2\delta_NI$, giving factor
$1+2\delta_N/(h_N-2\delta_N)\le(1-2\theta_N)^{-1}$.
This direct original-moment bound, unlike a still-to-be-selected
resolvent depth, supplies an unconditional finite entrywise
precision budget for any prescribed original four-endpoint error.

\section{Exact complex transport, inverse evaluation, and full fibres}

Retain the full original coefficient map, including all phases:
\begin{equation}\label{mprec:eq:transport}
 \begin{gathered}
 (T_N)_{rj}=\binom jr c^{j-r}i^r,\quad
 (T_N^{-1})_{rj}=\binom jr(-c)^{j-r}i^{-j}\quad(r\le j),\\
 L_N=\Pi_NT_N,\quad H_S=L_N^*H_NL_N,\quad
 \det T_N=i^{N(N+1)/2},\quad
 \det L_N=(\det\Pi_N)i^{N(N+1)/2},\\
 A=\Lambda J_{S,N}L_N^{-1}.
 \end{gathered}
\end{equation}
Other triangular entries vanish and $\Pi_N$ orders even powers
before odd powers. Expanding $(c+iy)^j$ proves the forward map;
expanding $y^j=i^{-j}(S-c)^j$ proves its inverse. The original
physical Taylor unit remains in $J_{S,N}$, and $A$ is onto the
same original observation image, not replaced by its norm.

For any positive source matrix $D$ in these parity coordinates,
\begin{equation}\label{mprec:eq:Q}
 Q_A(D)=(AD^{-1}A^*)^{-1},\quad
 \min_{Az=u}z^*Dz=u^*Q_A(D)u,\quad
 z_{\min}=D^{-1}A^*Q_A(D)u.
\end{equation}
The displayed minimizer is in the fibre and is $D$-orthogonal to
$\ker A$; subtraction and expansion of the square prove both
identities. Its original $S$ lift is $L_N^{-1}z_{\min}$.
For either outward source pair above, the same-fibre minimum gives
\begin{equation}\label{mprec:eq:Qbounds}
 Q^-:=Q_A(D^{\rm lo})\preceq Q_A(H_N)\preceq
 Q^+:=Q_A(D^{\rm up}).
\end{equation}
The source relative factor is preserved exactly. Any original
specified source subspace uses its same inclusion and restricted
fibre, never a silently enlarged polynomial source.

No further unverified inverse is hidden in \eqref{mprec:eq:Qbounds}.
Here is an explicit evaluation margin. For the actual onto $A$
with $d>0$ rows put
\begin{equation}\label{mprec:eq:Afloor}
 a_A=\frac{\det(AA^*)}{[\mprectr(AA^*)]^{d-1}}>0,
 \quad b_A=\mprectr(AA^*).
\end{equation}
The determinant/trace eigenvalue proof in Lemma \ref{mprec:lem:J}
gives $a_AI\preceq AA^*\preceq b_AI$. This is the literal
original map, including its jet and complex coefficients. If
$lI\preceq D\preceq UI$, then
\begin{equation}\label{mprec:eq:obsmargins}
 \frac{a_A}{U}I\preceq AD^{-1}A^*\preceq\frac{b_A}{l}I,
\qquad \frac l{b_A}I\preceq Q_A(D)\preceq\frac U{a_A}I.
\end{equation}
These margins have already proved native choices, so the
inequality is not an additional positivity premise. For the direct
moment endpoints take, simultaneously,
\[
 l=h_N-2\delta_N,\qquad
 U=\sum_{\epsilon=0}^1\sum_{i=0}^{n_\epsilon}
                 U_{\nu_\epsilon,2i}+2\delta_N.
\]
The exact native source has norm at most the displayed trace cap,
and each endpoint is within $2\delta_N$ of it. For the finite
resolvent endpoints take $l=g_N^*-2\eta_-$ (or its displayed
untilted replacement) and $U=\mprectr G^{\rm up}$, increased by any
certified arithmetic error in that trace. Positivity and the order
in \eqref{mprec:eq:sourcecert} prove these common endpoint bounds.
Thus a certified numerical center for $D^{-1}$ with error $r_D$
gives an observation-center error at most $b_Ar_D$, before adding
its own certified multiplication roundoff. Requiring that total
error $\delta_B\le a_A/(4U)$ gives the outward observation
inverse by \eqref{mprec:eq:inversebracket}, with $\ell=a_A/U$.
One can meet this stated precision using exact arithmetic or
\eqref{mprec:eq:residualinverse}. It is not inferred from a midpoint
determinant. If $A$ itself is evaluated numerically with error
$\delta_A$ and midpoint $\widehat A$, its additional product
error against a source inverse of norm at most $1/l$ is at most
$(2\mprecnorm{A}\delta_A+\delta_A^2)/l$ (or replace $\mprecnorm{A}$
by $\mprecnorm{\widehat A}+\delta_A$). Its original coefficient
enclosures are then separate inputs; this note does not manufacture
them. The determinant in \eqref{mprec:eq:Afloor} is an exact formula
in the original onto map, not a universal lower bound independent
of that map. The source-relative tolerance below does not require
its numerical evaluation.

In particular, if both the source inverse and the map are
approximated, the combined product radius, before multiplication
roundoff, is explicitly
\[
 \mprecnorm{AD^{-1}A^*-\widehat A R_0\widehat A^*}
 \le \frac{2\mprecnorm{A}\delta_A+\delta_A^2}{l}
                  +\mprecnorm{\widehat A}^{\,2}r_D.
\]
This follows by first changing $A$ against the exact inverse,
then changing the inverse against $\widehat A$. It retains the
mixed effect of simultaneous map and inverse errors rather than
incorrectly using the coefficient $b_A$ for both centers.

\section{Every complex cross term and every endpoint sign}

For any fixed original observation columns $U=(u_1,u_2,w)$,
the entire $3$ by $3$ Gram is bounded by
\[
 U^*Q^-U\preceq U^*Q_A(H_N)U\preceq U^*Q^+U.
\]
Its actual common remainder $E=U^*(Q_A(H_N)-Q^-)U$ satisfies
$0\preceq E\preceq U^*(Q^+-Q^-)U$; in particular
\begin{equation}\label{mprec:eq:cross}
 |u^*(Q_A(H_N)-Q^-)v|^2
 \le(u^*\Delta Q u)(v^*\Delta Q v),\qquad\Delta Q=Q^+-Q^-.
\end{equation}
This follows by Cauchy--Schwarz in the positive remainder. It
retains the single PSD matrix, including every principal minor and
the full determinant constraint, not three independently chosen
complex errors. For example, if $E_{ii}=\alpha_i$ and
$E_{ij}=z_{ij}$ for $i<j$, its constraints include $\alpha_i\ge0$,
$|z_{ij}|^2\le\alpha_i\alpha_j$, and the full determinant
\[
 \alpha_1\alpha_2\alpha_3
 +2\Re(z_{12}z_{23}\overline{z_{13}})
 -\alpha_1|z_{23}|^2-\alpha_2|z_{13}|^2-\alpha_3|z_{12}|^2
 \ge0.
\]
Expansion of the determinant proves this expression; the same
principal-minor constraints apply to $U^*\Delta Q U-E$.
In the original currents set
$u_i=\Lambda a_i$, $a_1=b_N$, $a_2=b_{N+1}$, and
$w=\Lambda v_\lambda$ with its original terminal class and
$\omega_N>0$. Put $B_i=u_i^*Q_A(H_N)w$, $B_i^-=u_i^*Q^-w$,
$\epsilon_i=B_i-B_i^-$ and
$e_i=[(u_i^*\Delta Q u_i)(w^*\Delta Q w)]^{1/2}$.
The complete signed product is
\begin{equation}\label{mprec:eq:current}
 \frac{2}{\omega_N}\Re(\overline B_1B_2)
 =\frac{2}{\omega_N}\Re\{
 \overline{B_1^-}B_2^-+\overline{B_1^-}\epsilon_2
 +\overline{\epsilon_1}B_2^-+\overline{\epsilon_1}\epsilon_2\}.
\end{equation}
Its difference from the first term has absolute value at most
$2(|B_1^-|e_2+|B_2^-|e_1+e_1e_2)/\omega_N$.
No cross term, conjugation, remainder, or phase is assigned zero.

For the actual image dimension $d_N$, define
$e_N=\log\det Q_N^+-\log\det Q_N^-\ge0$ in the unchanged
frame. A source relative factor $r_N\ge1$ gives
\begin{equation}\label{mprec:eq:logbound}
 0\le\log\det Q_N-\log\det Q_N^-\le e_N\le d_N\log r_N.
\end{equation}
The proof conjugates by $(Q_N^-)^{-1/2}$ and multiplies the
eigenvalues in $[1,r_N]$. This auxiliary computation does not
change the specified determinant or frame. The empty determinant
is one. At exactly the original window
$(N_1,N_2,N_3,N_4)=(q-1,q,2q-1,2q)$ let
\[
 F=\log\frac{\det Q_{N_1}\det Q_{N_2}}
                   {\det Q_{N_3}\det Q_{N_4}},\quad
 F^-=\log\frac{\det Q_{N_1}^-\det Q_{N_2}^-}
                     {\det Q_{N_3}^-\det Q_{N_4}^-}.
\]
The complete signed enclosure is
\begin{equation}\label{mprec:eq:four}
 \boxed{-e_{N_3}-e_{N_4}\le F-F^-
                         \le e_{N_1}+e_{N_2}.}
\end{equation}
Equivalently the lower ratio uses lower numerator matrices and
upper denominator matrices, and the upper ratio reverses these.

For a prescribed absolute four-endpoint tolerance $t_*>0$, the
direct actual-moment route supplies the explicit endpoint budgets
\begin{equation}\label{mprec:eq:finalbudget}
 \boxed{\theta_N=\min\{1/4,t_*/(8d_N)\},\qquad
 \varepsilon_N\le
 \frac{\theta_Nd_\nu\min(g_{n_0},g_{n_1})}
      {\max(n_0+1,n_1+1)}\quad(d_N>0).}
\end{equation}
Indeed \eqref{mprec:eq:directrelative} gives
$e_N\le-d_N\log(1-2\theta_N)\le4d_N\theta_N\le t_*/2$;
the middle inequality follows by integrating
$(1-u)^{-1}\le2$ over $0\le u\le2\theta_N\le1/2$.
If the minimum selects $1/4$, then $t_*\ge2d_N$ and the last
inequality still holds. Equation \eqref{mprec:eq:four} now proves
$|F-F^-|\le t_*$. Endpoints of rank zero have zero determinant
error and require no such budget. Taking the minimum of the four
positive budgets supplies one compatible shared moment precision.
No observation conditioning hypothesis is needed for this exact
relative result, although numerical evaluation of the endpoint
matrices still follows \eqref{mprec:eq:obsmargins}.

If numerical log determinants are returned, retain their outward
intervals too. Write $l_i^-\le\log\det Q_i^-\le u_i^-$ and
$l_i^+\le\log\det Q_i^+\le u_i^+$. The fully outward real
return interval is
\[
 [\,l_1^-+l_2^--u_3^+-u_4^+,\quad
       u_1^++u_2^+-l_3^--l_4^-\,].
\]
For an absolute error target $t_*$ on a reported approximation
to $F^-$, use \eqref{mprec:eq:finalbudget} with $t_*/2$ instead of
$t_*$ and certify each of the four numerical lower-endpoint log
values to absolute error at most $t_*/8$. The signed sum then
has numerical error at most $t_*/2$, and the total return error
is at most $t_*$. This separates, and includes, the actual moment
error and the finite arithmetic error.

\section{The same conductor and its unchanged signed-state kernel}

On the already established original conductor domain \cite{mprec:TR},
retain $C=\bar C\Lambda$ with $\bar C$ onto its actual image.
The exact conductor metric and outward endpoints are
\begin{equation}\label{mprec:eq:conductor}
 T_N=(\bar C Q_N^{-1}\bar C^*)^{-1},\quad
 T_N^\pm=(\bar C(Q_N^\pm)^{-1}\bar C^*)^{-1}
       =[(\bar CA)(D^\pm)^{-1}(\bar CA)^*]^{-1}.
\end{equation}
Substitution proves the last equality; minimizing on the full same
fibre $\bar Cu=v$ proves $T_N^-\preceq T_N\preceq T_N^+$
with the identical source-relative factor. This includes the entire
conductor kernel. Its numerical inverse floor is
\eqref{mprec:eq:obsmargins} with the literal map $\bar CA$.
Keep the original cubic inclusion $I_W$, coordinate map $\Psi$,
and signed-state map $B_{\mathcal J}=O(I-D_{\mathcal J})O^{-1}E$.
For the original label vector $\alpha_0$ the outward signed-state
energies are precisely
\[
 \alpha_0^*B_{\mathcal J}^*\Psi^*I_W^*T_N^-I_W\Psi
               B_{\mathcal J}\alpha_0
 \le\alpha_0^*B_{\mathcal J}^*\Psi^*I_W^*T_NI_W\Psi
               B_{\mathcal J}\alpha_0
 \le\alpha_0^*B_{\mathcal J}^*\Psi^*I_W^*T_N^+I_W\Psi
               B_{\mathcal J}\alpha_0.
\]
These are literal congruences of \eqref{mprec:eq:conductor}. The
terminal-class residual of the original EQ20--21 is retained,
not set to zero by these positive energy bounds. Its remaining
complex pairings use \eqref{mprec:eq:cross} with the original columns.
Conductor determinants, or a specified positive definite
restriction, use \eqref{mprec:eq:four} and \eqref{mprec:eq:finalbudget} with
their actual ranks, frames and the same four endpoints.

\section{Finite native scalar inputs and what the precision costs}

The rational numbers $g_t$, all pole values, determinant phases,
dimension factors and error coefficients above are fully specified
finite expressions. No native integral has been numerically
evaluated without supplied values of the original quartet data.
The remaining scalar inputs are precisely the original
$\vartheta_h$, $M_\alpha$, TW7/BT12 constants, and finitely many
actual moments. Downward bounds for $a_k,d_\nu$ and upward bounds
for $A_k,U_{\lambda,r}$ may replace their exact values in every
budget. This preserves all inequalities; it never silently treats
an uncertified midpoint constant as a lower bound.

For the untilted route the moments can be reduced further to
original onefold scalar integrals. Write
$\eta_l=\int y^lw_h(y)dy$ and
$\zeta_L=\int y^Lm_k(y)dy$. Evenness gives $\eta_{2j+1}=0$,
and Tonelli followed by the finite multinomial formula gives
\begin{equation}\label{mprec:eq:multinomial}
 \int x^r d\nu_\epsilon=\zeta_{2(r+\epsilon)},\quad
 \zeta_L=\sum_{l_1+\cdots+l_k=L}
 \frac{L!}{l_1!\cdots l_k!}\prod_{j=1}^k\eta_{l_j}.
\end{equation}
Odd summands vanish, and all remaining factors are nonnegative.
Every monomial product in the finite expansion is absolutely
integrable by the original exponential moment. Fubini therefore
justifies the factorization even for the terms containing odd
powers, before evenness makes those terms zero.
For $L=0$ this is the actual mass $\mu_h^k$; it is not one.
The original exponential moment gives the explicit caps
\begin{equation}\label{mprec:eq:etaone}
 0\le\eta_{2j}\le(2j)!\alpha^{-2j}M_\alpha,
\end{equation}
because $\cosh(\alpha y)\ge(\alpha y)^{2j}/(2j)!$ and
$\int\cosh(\alpha y)w_h=M_\alpha$.
Suppose original moments through $\zeta_{2R}$ are wanted at
absolute error $\varepsilon$. Put
$U=\max_{0\le j\le R}(2j)!\alpha^{-2j}M_\alpha$.
If every even onefold input is enclosed with midpoint error
\begin{equation}\label{mprec:eq:onefoldbudget}
 \tau\le\min\left\{1,
 \frac{\varepsilon}{k^{2R+1}(U+1)^{k-1}}\right\},
\end{equation}
then the finite polynomial \eqref{mprec:eq:multinomial}, evaluated
exactly on those midpoints, has error at most $\varepsilon$
for every even degree through $2R$. To prove it, telescope each
$k$-fold product; its error is at most
$k\tau(U+1)^{k-1}$. The sum of all multinomial coefficients
at degree $L$ is $k^L$, and the even-only subsum is smaller.
For $L\le2R$ their product is bounded by the denominator in
\eqref{mprec:eq:onefoldbudget}. Thus the four-endpoint budget is an
explicit finite accuracy request on original onefold integrals,
not on an unidentified proxy measure.

Their infinite tails also have an original explicit certificate.
For an even degree $L$ and $Y\ge\max\{1,L/\alpha\}$,
\begin{equation}\label{mprec:eq:tail}
 \int_{|y|>Y}|y|^Lw_h(y)dy
 \le2M_\alpha Y^Le^{-\alpha Y},\quad
 \int_{|y|>Y}|y|^Lm_k(y)dy
 \le2M_\alpha^kY^Le^{-\alpha Y}.
\end{equation}
Indeed $y^Le^{-\alpha y}$ decreases for $y\ge L/\alpha$,
and evenness gives $\int e^{\alpha|y|}w_h\le2M_\alpha$.
Tonelli proves $\int e^{\alpha y}m_k=M_\alpha^k$, giving
the second bound. A tilted moment tail has the additional factor
at most $\alpha$. Repeatedly doubling $Y$ from the stated lower
threshold attains any prescribed positive tail allowance, since
$Y^Le^{-\alpha Y}\to0$. The remaining compact integrals of
the original $w_h$, or of $s(y^2)m_k(y)$ if one chooses the
tilted route, still require certified evaluation. Their numerical
values are not supplied by these tail inequalities or by the
moment-matrix calculation.

For absolute input error $2^{-p}$, \eqref{mprec:eq:uniformbudget}
requires the explicit integer
\[
 p\ge\max\left\{0,
 \left\lceil\log_2\frac{t+1}{\theta d_\lambda g_t}\right\rceil
 \right\}.
\]
There is no hidden unknown eigenvalue. Formula \eqref{mprec:eq:gsize}
shows an elementary sufficient cost with
$-\log_2g_t\le4t^2+2t+t\log_2(8/3)
 +(t+1)\log_2(2t+1)$.
These unaltered monomial-coordinate floors are deliberately
conservative and may demand very high precision at the original
large degrees. This proof establishes finite, explicit tolerances
and outward certificates, not practical low-cost evaluation of
the native integrals or a new phase/asymptotic conclusion.

The companion exact tests verify rational interval Grams,
determinants, all polynomial weights, moment perturbations, inverse
brackets, the deflated Radau margin, complex rank-one cross terms,
original triangular maps and masses, quotient/conductor fibres,
and all four determinant signs. Test measures serve only as exact
checks of the proved universal identities. They do not replace the
native measures in any statement above.

\begin{thebibliography}{9}
\bibitem{mprec:ZDS} J\"orn Zimmerling, Vladimir Druskin and Valeria Simoncini,
\emph{Monotonicity, bounds and acceleration of Block Gauss and
Gauss--Radau quadrature for computing $B^T\phi(A)B$}, original
author source of arXiv:2407.21505v3, 9 January 2025, Sections 3--5.
\url{https://arxiv.org/abs/2407.21505v3}.
The abstract-page title uses ``extrapolation''; the title here
identifies the retained version-three author source. Its complete
source was read by the integrating task. This note uses it for
human mathematical provenance, and proves its native conditioning
and perturbation claims independently.
\bibitem{mprec:CAI} Split-Zero programme, \emph{The complete original action
relative to its arithmetic total}, complete provider
\path{CAI_COMPLETE_PREREQUISITE.tex}, CAI16--17 and its original
TW7/BT12 constant definitions. The exact displayed envelope and
constants are retained above.
\bibitem{mprec:HG} Split-Zero programme, \emph{Subexponential heat cost and
the original arithmetic currents}, complete provider
\path{HEAT_GROWTH_AND_ORIGINAL_CURRENTS.tex}, HG21, arbitrary
fixed primary multiplicity and the full-measure comparison.
\bibitem{mprec:PP} Split-Zero programme, \emph{Exact native parity transport
into the Pollaczek pair and certified finite source--observation
bounds}, complete proof
\path{independent_parity_pollaczek/PARITY_POLLACZEK_NATIVE_RECEIVER.tex},
20 September 2026. The unscaled Pollaczek ratio there retains the
human source A. I. Aptekarev, G. L\'opez Lagomasino and
A. Mart\'{\i}nez--Finkelshtein, \emph{Strong asymptotics for the
Pollaczek multiple orthogonal polynomials ensembles},
arXiv:1410.1261v1 (2014). No asymptotic or replacement weight from
that paper is imported here.
\bibitem{mprec:FM} Split-Zero programme, \emph{Certified finite native moment
enclosures for the parity resolvents}, complete proof
\path{independent_resolvent_moments/NATIVE_RESOLVENT_MOMENT_ENCLOSURES.tex},
20 September 2026. Read completely as a prerequisite.
\bibitem{mprec:RR} Split-Zero programme, \emph{One scalar resolvent
uncertainty in the original polynomial observation}, complete proof
\path{NATIVE_RESOLVENT_RANK_ONE.tex}, RR1--16, 20 September 2026.
Read completely as a prerequisite; its full mass, deflation,
rank-one phase, original maps and endpoint signs are retained.
\bibitem{mprec:TR} Split-Zero programme, \emph{Native resolvent truncation at
growing degree and the original conductor minimum}, complete proof
\path{NATIVE_TRUNCATION_CONDUCTOR_RETURN.tex}, TR1--10,
20 September 2026, retaining the complete original EQ conductor
and signed-state maps, including their terminal residuals.
\end{thebibliography}

\endgroup

\clearpage\begingroup
\part*{Complete Burnol-to-original-theta transfer}
\newtheorem{mburntheorem}{Theorem}[section]
\newtheorem{mburnproposition}[mburntheorem]{Proposition}
\newtheorem{mburnlemma}[mburntheorem]{Lemma}
\newtheorem{mburncorollary}[mburntheorem]{Corollary}
\newcommand{\mburnC}{\mathbb C}\newcommand{\mburnR}{\mathbb R}
\newcommand{\mburnB}{\mathscr B}\newcommand{\mburnHH}{\mathcal H}
\newcommand{\mburnML}{\mathcal M_{\!L}}\newcommand{\mburnMR}{\mathcal M_{\!R}}
\newcommand{\mburnPP}{\mathcal P}\newcommand{\mburnFc}{\mathcal F_+}
\newcommand{\mburnnorm}[1]{\left\lVert#1\right\rVert}


\section{Fixed spaces, transforms, and the exact human input}

All vector spaces are complex. The physical coordinate and Hilbert
measure throughout are $x>0$ and $dx$. Put
\begin{equation}\label{mburn:eq:objects}
\begin{gathered}
 \mburnHH=L^2((0,\infty),dx),\quad D=-x\partial_x,\quad
 IF(x)=x^{-1}F(x^{-1}),\\
 V=\{\phi\in\mathcal S(\mburnR):\phi(-x)=\phi(x),\ 
       \phi(0)=0,\ \int_\mburnR\phi(x)\,dx=0\},\\
 \Theta\phi(x)=2\sum_{n\geq1}\phi(nx),\quad
 p_{b,j}(F)=\sup_{x>0}x^b|D^jF(x)|,\\
 \mburnB=\{F\in C^\infty(0,\infty):p_{b,j}(F)<\infty
                          \ (b\in\mathbb Z,j\geq0)\},
 \quad Q_\Theta=\mburnB/\Theta V.
\end{gathered}
\end{equation}
The topology of $V$ is its Schwartz subspace topology; the topology
of $\mburnB$ is generated by every displayed $p_{b,j}$. The map to the
quotient is denoted $q$. Inner products in this note are
$\langle f,g\rangle=\int\overline f g\,dx$; Burnol's evaluators will
retain instead his stated complex bilinear pairing
$[f,g]=\int fg\,dx$.

The two Mellin transforms are
\begin{equation}\label{mburn:eq:mellin}
 \mburnML F(s)=\int_0^\infty F(x)x^{s-1}\,dx,\qquad
 \mburnMR f(s)=\int_0^\infty f(x)x^{-s}\,dx.
\end{equation}
On $\Re s=1/2$ these are unitary transforms from $\mburnHH$ to
$L^2(\mburnR,dt/(2\pi))$, with $s=1/2+it$, interpreted by $L^2$
limits. Indeed $x=e^y$ sends $f$ to $e^{y/2}f(e^y)$ with the same
norm, and the transforms become its Fourier transforms with the
respective positive and negative exponentials. In particular
\begin{equation}\label{mburn:eq:parseval}
 \mburnML(If)=\mburnMR f,\qquad
 \langle F,H\rangle=\int_\mburnR
       \overline{\mburnML F(1/2+it)}\mburnML H(1/2+it)\,\frac{dt}{2\pi}.
\end{equation}

The cosine transform is the unitary involution
\[
 \mburnFc f(x)=2\int_0^\infty\cos(2\pi xy)f(y)\,dy
\]
initially on integrable square-integrable functions and then by
Plancherel. Define the closed subspace
\[
 L_1=\{f\in\mburnHH:f\text{ is constant a.e. on }(0,1),\quad
                      \mburnFc f\text{ is constant a.e. on }(0,1)\}.
\]
It is closed because the constant functions form a closed subspace
of $L^2(0,1)$ and restriction and $\mburnFc$ are bounded. Denote its
Hilbert norm by the original $dx$ norm, not a new measure.

Here is precisely the human input from Burnol \cite{mburn:Burnol}.
Sections 2--6 of the
\href{https://arxiv.org/abs/math/0203120v7}{versioned original author source}
are used.
For $f\in L_1$, with $G=\mburnMR f$ and
\begin{equation}\label{mburn:eq:abg}
 b(s)=\pi^{-s/2}\Gamma(s/2),\qquad
 a(s)=s(s-1)b(s),\qquad g(s)=a(s)\zeta(s)=2\xi(s),
\end{equation}
$bG$ continues meromorphically to the plane, with possible poles
only at $0,1$. At any $w\ne0,1$ its derivatives are continuous
linear functionals on $L_1$, represented by $Y^1_{w,k}$ as
\begin{equation}\label{mburn:eq:burnol-eval}
 [f,Y^1_{w,k}]=(bG)^{(k)}(w).
\end{equation}
For each nontrivial zero $\rho$ of exact multiplicity $m_\rho$,
there is $X_{\rho,\ell}\in L_1$ with
\begin{equation}\label{mburn:eq:X}
 \mburnMR X_{\rho,\ell}(s)=\frac{\zeta(s)}{(s-\rho)^\ell},
                 \qquad 1\leq\ell\leq m_\rho.
\end{equation}
Both families $\{X_{\rho,\ell}\}$ and
$\{Y^1_{\rho,k}:0\leq k<m_\rho\}$ are complete and minimal in
$L_1$. These are respectively Proposition \texttt{thm:mellinLa},
Proposition \texttt{zetaprop}, Theorem \texttt{thmB}, and its
Section 6 proof in the author source. These established completeness
theorems are the imported results; no zero-location assertion is
imported. All maps, coefficients, topological estimates and metric
identities below are proved here.

\section{The full-factor integral transform, including strong topology}

Put
\begin{equation}\label{mburn:eq:kernel}
 \phi_*(x)=(4\pi^2x^4-6\pi x^2)e^{-\pi x^2},\qquad K(x)=2\phi_*(x).
\end{equation}
These are the original functions without a scale change. For
$\Re s>-2$, direct Gaussian integration gives
\begin{align}
 \mburnML K(s)
 &=\pi^{-s/2}\{4\Gamma(s/2+2)-6\Gamma(s/2+1)\}\nonumber\\
 &=s(s-1)\pi^{-s/2}\Gamma(s/2)=a(s).\label{mburn:eq:MK}
\end{align}
The recurrence $\Gamma(z+1)=z\Gamma(z)$ proves the second equality,
including its continuation through the removable value at $s=0$.
Both $K(0)=0$ and $\int_0^\infty K=\mburnML K(1)=0$ hold. Thus
$\phi_*\in V$. Moreover, with the ordinary Fourier kernel
$e^{-2\pi i x\xi}$, $\widehat{D\phi}=(1-D)\widehat\phi$.
Since the Gaussian is Fourier invariant and
$D(D-1)e^{-\pi x^2}=\phi_*$, this identity proves
\begin{equation}\label{mburn:eq:selfdual}
 \mburnFc K=K.
\end{equation}

\begin{mburntheorem}[The exact transfer on its full domains]\label{mburn:thm:T}
The formula
\begin{equation}\label{mburn:eq:T}
 (Tf)(x)=\int_0^\infty K(xt)f(t)\,dt
       =(K*_\times If)(x),\qquad
 (u*_\times v)(x)=\int_0^\infty u(x/y)v(y)\,\frac{dy}{y},
\end{equation}
defines a bounded injective map $T:\mburnHH\to\mburnHH$, with
\begin{equation}\label{mburn:eq:Tsymbol}
 \mburnML(Tf)(s)=a(s)\mburnMR f(s)\quad(\Re s=1/2),\qquad
 \mburnnorm T_{\mburnHH\to\mburnHH}=\sup_{t\in\mburnR}|a(1/2+it)|.
\end{equation}
Its restriction is a continuous linear injection
$T:L_1\to\mburnB$. It satisfies the exact reflection identity
\begin{equation}\label{mburn:eq:Tsym}
 T\mburnFc=IT.
\end{equation}
For an integer $M\geq2$ define finite constants
\begin{equation}\label{mburn:eq:constants}
 k_{M,j}=\sup_{u\geq1}u^M|D^jK(u)|,\qquad
 C_{M,j}=k_{M,j}\left(\frac1{M-1}+\frac1{\sqrt{2M-1}}\right).
\end{equation}
For every $f\in L_1$, every $b\in\mathbb Z$ and $j\geq0$,
with $M=\max(2,b,1-b)$,
\begin{equation}\label{mburn:eq:strongbound}
 p_{b,j}(Tf)\leq
 \left(C_{M,j}+\sum_{r=0}^j\binom jr C_{M,r}\right)\mburnnorm f_2.
\end{equation}
For $f\in L_1$ the equality $\mburnML(Tf)=a\mburnMR f$ holds on the
whole plane by continuation, and its left side is entire.
\end{mburntheorem}
\begin{proof}
For each $x>0$, Cauchy--Schwarz gives absolute convergence of
\eqref{mburn:eq:T} and the bound
$|Tf(x)|\leq x^{-1/2}\mburnnorm K_2\mburnnorm f_2$.
The change $y=1/t$ proves the convolution formula, with no extra
factor. On compact positive $x$ intervals the kernels and their
ordinary derivatives vary continuously in $L^2(dt)$; the same
Gaussian majorants dominate their difference quotients. Consequently
$Tf$ is smooth and $D^jTf(x)=\int(D^jK)(xt)f(t)\,dt$.

The multiplicative convolution acts on $\mburnHH$ with norm at most
$\int_0^\infty|K(y)|y^{-1/2}\,dy<\infty$:
the norm of $v(x/y)$ is $y^{1/2}\mburnnorm v_2$, so the triangle
inequality for the Bochner integral gives precisely this bound.
For compactly supported bounded $f$ the Mellin multiplier identity
follows by absolute integration after taking Mellin transforms on
the critical line in the Fourier sense. Alternatively the above
logarithmic variable turns this weighted convolution into ordinary
$L^1*L^2$ convolution, for which Fourier transformation multiplies.
Density proves the identity for all $f\in\mburnHH$. Its multiplier
norm is the essential supremum of $|a|$. The function $a$ is
continuous and bounded on the critical line, tends to zero at both
ends by the vertical-line Gamma estimate, and is nonzero there
because $s(s-1)\ne0$ and Gamma has no zeros or poles there.
Thus essential supremum equals supremum, the exact multiplier
norm holds, and Mellin unitarity proves injectivity.

The cosine transform preserves the bilinear pairing on $\mburnHH$;
this follows first from its real symmetric kernel and then from
density. The cosine transform of $t\mapsto K(xt)$ is
$t\mapsto x^{-1}K(t/x)$ by \eqref{mburn:eq:selfdual} and substitution.
Applying the pairing identity gives $T\mburnFc f(x)=x^{-1}Tf(1/x)$,
which proves \eqref{mburn:eq:Tsym} on all of $\mburnHH$.

Suppose now $f\in L_1$ and $f=c$ a.e. on $(0,1)$. Then
$|c|\leq\mburnnorm f_2$. Integration by parts, with vanishing
endpoints, gives $\int D^jK=\int K=0$ for every $j$. For $x\geq1$
we therefore have the exact cancellation formula
\[
 D^jTf(x)=-\frac c x\int_x^\infty D^jK(u)\,du
                      +\int_1^\infty(D^jK)(xt)f(t)\,dt.
\]
The first integral is bounded by $k_{M,j}x^{1-M}/(M-1)$.
The second term is bounded by
\[
 \mburnnorm f_2x^{-1/2}\left(\int_x^\infty|D^jK(u)|^2du\right)^{1/2}.
\]
Combining these bounds proves
\begin{equation}\label{mburn:eq:tail}
 |D^jTf(x)|\leq C_{M,j}\mburnnorm f_2 x^{-M}\quad(x\geq1).
\end{equation}
The identity $D^jI=I(1-D)^j$ follows by the chain rule. Apply
\eqref{mburn:eq:Tsym} to $\mburnFc f$, also in $L_1$ with the same norm,
to obtain for $0<x\leq1$
\[
 |D^jTf(x)|\leq\mburnnorm f_2 x^{M-1}
                           \sum_{r=0}^j\binom jr C_{M,r}.
\]
The inequalities $M\geq b$ and $M\geq1-b$ prove
\eqref{mburn:eq:strongbound}. Thus $T:L_1\to\mburnB$ is continuous in
every original seminorm. The Mellin transform of an element of
$\mburnB$ is entire by endpoint domination; identity on the critical
line and uniqueness of meromorphic continuation give the last
assertion.
\end{proof}

The same cancellation is unavailable on unrestricted $\mburnHH$.
For example $f=\mathbf1_{(0,1)}$ gives
$Tf(x)=x^{-1}\int_0^xK(u)\,du$. Since
$K(u)=-12\pi u^2+O(u^4)$, this equals
$-4\pi x^2+O(x^4)$ as $x\to0$ and is not in $\mburnB$.
Thus the domain $L_1$ in the strong assertion is material.

\begin{mburnproposition}[Compactness and the exact quotient injection]
\label{mburn:prop:quotient}
The image of a bounded set of $L_1$ under $T$ is relatively
compact in $\mburnB$. In particular $T:L_1\to\mburnHH$ is compact and
has no bounded inverse on its image. Also
\begin{equation}\label{mburn:eq:intersection}
 T(L_1)\cap\Theta V=\{0\},\qquad
 qT:L_1\longrightarrow Q_\Theta\ \text{is continuous and injective}.
\end{equation}
Neither $T$ nor $qT$ is a topological embedding of the Hilbert
topology of $L_1$ into the indicated original strong topology.
\end{mburnproposition}
\begin{proof}
For any fixed $b,j$, use \eqref{mburn:eq:tail} and its reflected version
with $M>\max(b,1-b)$. The weighted derivative is then uniformly
small outside a sufficiently large compact positive interval for
all $f$ in a fixed Hilbert ball. On that interval the next
derivative bounds give uniform equicontinuity. Arzel\`a--Ascoli,
a diagonal subsequence over all compact intervals and derivative
orders, and the uniform tails give convergence in all $p_{b,j}$.
The limits of successive derivatives agree by integration on
compact intervals. Thus the limit belongs to $\mburnB$, proving
relative compactness directly. The inclusion $\mburnB\to\mburnHH$ is
continuous because
$\mburnnorm F_2^2\leq p_{0,0}(F)^2+p_{1,0}(F)^2$.

Every $\mburnML\Theta\phi=2\zeta\mburnML\phi$ vanishes to the complete
multiplicity at each nontrivial zero: first prove the equality in
$\Re s>1$ by absolute integration, then continue to $\Re s>-2$,
where $\mburnML\phi$ is holomorphic since $\phi(x)=O(x^2)$ at zero.
If $Tf=\Theta\phi$, all jets of $aG=s(s-1)bG$ at those zeros
vanish. The multiplier $s(s-1)$ is a nonzero local unit, so all
derivatives $(bG)^{(k)}(\rho)$ for $k<m_\rho$ vanish. The
completeness of the Burnol evaluators, in their bilinear pairing,
forces $f=0$: $[f,Y]=0$ on a dense span implies $[f,v]=0$ for
all $v\in L_1$, and one may take $v=\overline f\in L_1$.
This proves \eqref{mburn:eq:intersection}; continuity follows from
Theorem~\ref{mburn:thm:T} and the quotient topology.

The space $L_1$ is infinite dimensional (the distinct rational
functions in \eqref{mburn:eq:X} are linearly independent). Take an
orthonormal sequence $f_n$ in it. It converges weakly to zero, so
every $D^jTf_n(x)$ tends to zero by its $L^2$ kernel pairing.
Relative compactness just proved forces $Tf_n\to0$ in $\mburnB$:
any subsequence not converging to zero would have a convergent
subsequence whose pointwise limit must be zero. Then also
$qTf_n\to0$, while $\mburnnorm{f_n}_2=1$. This disproves inverse
continuity for each of the claimed embeddings and a bounded
Hilbert inverse.
\end{proof}

The strong closed-range theorem $\overline{\Theta V}^{\mburnB}=\Theta V$
for this exact source and target is proved in \cite{mburn:ClosedTheta}.
It makes $Q_\Theta$ Hausdorff. Proposition~\ref{mburn:prop:quotient}
does not replace that topology by a Hilbert quotient. In particular
the vanishing of a Hilbert quotient seminorm is not the vanishing
of the injectively received class in \eqref{mburn:eq:intersection}.

\section{Every finite divisor and every polynomial degree}

Let $Z$ be a nonempty finite set of actual nontrivial zeros, and
retain all their exact multiplicities:
\begin{equation}\label{mburn:eq:h}
 h(s)=\prod_{\rho\in Z}(s-\rho)^{m_\rho},\quad d=\deg h,
 \quad h_\rho(s)=h(s)/(s-\rho)^{m_\rho},\quad
 E_h=\prod_{\rho\in Z}\mburnC[t_\rho]/(t_\rho^{m_\rho}).
\end{equation}
Here $t_\rho=s-\rho$ is literal, and $j_hH$ records the divided
Taylor coefficients $H^{(j)}(\rho)/j!$. No symmetry of $Z$ is
required. The original physical unit and its inverse are
\begin{equation}\label{mburn:eq:unit}
 U_h=j_h(g/h),\qquad \varepsilon_h=U_h^{-1}=j_h(h/g).
\end{equation}
Both notations use local germs at $Z$; $g/h$ is entire and is
nonzero at every selected root because its full multiplicity
has been removed.

\begin{mburnlemma}[The original physical divided function]\label{mburn:lem:Fh}
There is a unique $F_h\in\mburnB$ with $\mburnML F_h=g/h$.
For $F\in\mburnB$ and $\mburnML F(\rho)=0$, the literal inverse is
\[
 S_\rho F(x)=x^{-\rho}\int_x^\infty F(t)t^{\rho-1}\,dt,
 \quad(D-\rho)S_\rho F=F,\quad
 \mburnML S_\rho F(s)=\frac{\mburnML F(s)}{s-\rho}.
\]
Thus $F_h$ is obtained by applying the $S_\rho$'s through
the complete multiset of roots to $f_0=\Theta\phi_*$.
\end{mburnlemma}
\begin{proof}
Poisson summation, with both endpoint moments of $\phi_*$ zero,
shows $f_0\in\mburnB$; all derivatives have the same rapid endpoint
bounds. Equations \eqref{mburn:eq:MK} and termwise Mellin integration
give $\mburnML f_0=\zeta a=g$. For $S_\rho$, rapid decrease bounds
the tail at infinity. The zero moment rewrites its defining
integral as $-x^{-\rho}\int_0^xF(t)t^{\rho-1}dt$, proving rapid
decrease at zero. Differentiation gives the displayed differential
identity and all derivative bounds. Mellin integration by parts
then gives the transform identity. At each iteration the required
zero persists until its last assigned factor is removed.
Mellin uniqueness on the critical line proves uniqueness and
independence of the division order.
\end{proof}

For every polynomial $P$, use literal monic Euclidean division
\begin{equation}\label{mburn:eq:division}
 P=hQ+R,\qquad \deg R<d.
\end{equation}
The proper fraction has coefficients
\begin{equation}\label{mburn:eq:partial}
 \frac{R(s)}{h(s)}=
 \sum_{\rho\in Z}\sum_{\ell=1}^{m_\rho}
        \frac{c_{\rho,\ell}(R)}{(s-\rho)^\ell},\qquad
 c_{\rho,\ell}(R)=
 \frac{1}{(m_\rho-\ell)!}
 \left(\frac d{ds}\right)^{m_\rho-\ell}
                  \left(\frac R{h_\rho}\right)(\rho).
\end{equation}
This coefficient formula includes all other-root factors. More
explicitly, with $H_\rho=\prod_{\sigma\ne\rho}
(\rho-\sigma)^{m_\sigma}$, let
\begin{equation}\label{mburn:eq:other-roots}
 \alpha_{\rho,n}=H_\rho^{-1}(-1)^n
 \sum_{\substack{n_\sigma\geq0\ (\sigma\ne\rho)\\
                         \sum_{\sigma\ne\rho}n_\sigma=n}}
 \prod_{\sigma\ne\rho}
  \binom{m_\sigma+n_\sigma-1}{n_\sigma}
                    (\rho-\sigma)^{-n_\sigma}.
\end{equation}
The empty-product convention gives $\alpha_{\rho,0}=1$ and
$\alpha_{\rho,n}=0$ for $n>0$ when $Z=\{\rho\}$.
Then $1/h_\rho(\rho+t)=\sum_{n\geq0}\alpha_{\rho,n}t^n$ and
\[
 c_{\rho,\ell}(R)=\sum_{j=0}^{m_\rho-\ell}
       \frac{R^{(j)}(\rho)}{j!}\alpha_{\rho,m_\rho-\ell-j}.
\]
The binomial series of each factor of $1/h_\rho$ proves
\eqref{mburn:eq:other-roots}. Taylor expansion proves
\eqref{mburn:eq:partial}; subtracting all its principal parts leaves
a rational function without poles tending to zero at infinity,
so the equality is global, including repeated roots.

\begin{mburntheorem}[All-degree transfer, with the original relation term]
\label{mburn:thm:all-degree}
Define
\begin{equation}\label{mburn:eq:fR}
 f_{R/h}=\sum_{\rho,\ell}c_{\rho,\ell}(R)X_{\rho,\ell}\in L_1.
\end{equation}
For every polynomial $P$, with \eqref{mburn:eq:division},
\begin{equation}\label{mburn:eq:all-degree}
 \boxed{\quad P(D)F_h=T f_{R/h}+\Theta\bigl(Q(D)\phi_*\bigr).\quad}
\end{equation}
This is equality in $\mburnB$ with its original topology and hence
also equality of functions and of $dx$ Hilbert vectors. Every
full zero jet and the original quotient class are preserved.
The direct expression $\zeta P/h$ is the right Mellin transform
of an element of $L_1$ if and only if $Q=0$.
\end{mburntheorem}
\begin{proof}
Membership in $L_1$ follows from \eqref{mburn:eq:X} and the finite
linear combination. Its Mellin transform is $\zeta R/h$ by
\eqref{mburn:eq:partial}. Theorem~\ref{mburn:thm:T} gives
$\mburnML T f_{R/h}=a\zeta R/h=gR/h$. The polynomial $Q(D)$
preserves $V$: $D$ preserves evenness and value zero at zero,
and integration by parts gives $\int D\phi=\int\phi=0$.
Its theta image has transform $gQ$. The transform of the left
side is $Pg/h$. Identity \eqref{mburn:eq:division} and Mellin
uniqueness prove \eqref{mburn:eq:all-degree}, including its exact
source term. This term has every original zero jet equal to zero.

If $\zeta P/h$ were a Burnol transform, subtract
$\mburnMR f_{R/h}$ to obtain a vector of $L_1$ with transform
$\zeta Q$. All its completed nontrivial-zero jets vanish,
so evaluator completeness forces that vector to be zero, hence
$\zeta Q=0$ and $Q=0$. The reverse implication has just been
proved. Thus omission of the polynomial-source term is not
an admissible all-degree transfer; the exact direct-sum
decomposition is \eqref{mburn:eq:all-degree}.
\end{proof}

For $u\in E_h$ let $r_h(u)$ be the unique polynomial of degree
less than $d$ with $j_hr_h(u)=\varepsilon_hu$. The original
section is consequently
\begin{equation}\label{mburn:eq:Rh}
 R_hu=r_h(u)(D)F_h=T f_{r_h(u)/h},\qquad
 j_h\mburnML R_hu=u.
\end{equation}
It has zero full jets at every nontrivial zero outside $Z$.
If $Z\subset Z'$ and $u'$ extends $u$ by zero, then
$r_{h'}(u')=(h'/h)r_h(u)$, as follows by its jets and its
degree $<\deg h'$. Thus both sides of \eqref{mburn:eq:Rh} agree
as actual functions under finite-divisor enlargement.

\section{Exact confluent evaluation and both biorthogonal systems}

Fix a zero $\rho$ of multiplicity $m$. Use the convergent local
expansions
\begin{equation}\label{mburn:eq:local}
 \frac{\zeta(\rho+t)}{t^m}=z_\rho(t)
  =\sum_{n\geq0}\frac{\zeta^{(m+n)}(\rho)}{(m+n)!}t^n,
 \quad b(\rho+t)=\sum_{n\geq0}\beta_{\rho,n}t^n,
 \quad d_\rho(t)=b(\rho+t)z_\rho(t)=\sum_{n\geq0}d_{\rho,n}t^n.
\end{equation}
Here none of the Gamma derivatives is omitted:
\begin{equation}\label{mburn:eq:gamma-coeff}
 \beta_{\rho,n}=\pi^{-\rho/2}
 \sum_{v=0}^n
 \frac{(-\log\pi/2)^{n-v}\Gamma^{(v)}(\rho/2)}
                   {(n-v)!v!2^v},\qquad
 d_{\rho,n}=\sum_{v=0}^n\beta_{\rho,v}
                  \frac{\zeta^{(m+n-v)}(\rho)}{(m+n-v)!}.
\end{equation}
The coefficient $d_{\rho,0}=b(\rho)\zeta^{(m)}(\rho)/m!$
is nonzero. Define its reciprocal coefficients by
\begin{equation}\label{mburn:eq:reciprocal}
 \frac1{d_\rho(t)}=\sum_{n\geq0}c_{\rho,n}t^n,\qquad
 c_{\rho,0}=d_{\rho,0}^{-1},\quad
 c_{\rho,n}=-d_{\rho,0}^{-1}
                    \sum_{v=1}^nd_{\rho,v}c_{\rho,n-v}.
\end{equation}
Every coefficient with a negative subscript is defined to be zero.
The symbols $c_{\rho,n}$ in \eqref{mburn:eq:reciprocal} have one
nonnegative index; they are distinct from the explicitly
two-index partial-fraction coefficients $c_{\rho,\ell}(R)$.

\begin{mburntheorem}[The complete finite inverse, including factorials]
\label{mburn:thm:jets}
With row $k=0,\ldots,m-1$ and column $\ell=1,\ldots,m$, set
\begin{equation}\label{mburn:eq:M}
 M^\rho_{k\ell}=[X_{\rho,\ell},Y^1_{\rho,k}]
        =k!d_{\rho,k-m+\ell}.
\end{equation}
Cross-root pairings are zero. The exact inverse and determinant are
\begin{equation}\label{mburn:eq:Minv}
 (M^\rho)^{-1}_{\ell k}=\frac{c_{\rho,m-\ell-k}}{k!},\qquad
 \det M^\rho=(-1)^{m(m-1)/2}
              \left(\prod_{k=0}^{m-1}k!\right)d_{\rho,0}^{\,m}.
\end{equation}
In Burnol's bilinear convention the exact duals are
\begin{equation}\label{mburn:eq:duals}
 U_{\rho,k}=\sum_{\ell=1}^{m}(M^\rho)^{-1}_{\ell k}X_{\rho,\ell},
 \qquad
 Z_{\rho,\ell}=\sum_{k=0}^{m-1}(M^\rho)^{-1}_{\ell k}Y^1_{\rho,k},
\end{equation}
so $[U_{\rho,k},Y^1_{\sigma,j}]=\delta_{\rho\sigma}\delta_{kj}$
and $[X_{\sigma,j},Z_{\rho,\ell}]=\delta_{\rho\sigma}\delta_{j\ell}$.
For the Hermitian convention of this note the representing vectors
of these linear functionals are respectively
$\overline{Y^1_{\rho,k}}$ and $\overline{Z_{\rho,\ell}}$;
no unnoticed complex conjugation is inserted in \eqref{mburn:eq:duals}.
\end{mburntheorem}
\begin{proof}
At $\rho$ the completed transform of $X_{\rho,\ell}$ is
$t^{m-\ell}d_\rho(t)$. Its $k$th derivative is
$k!d_{\rho,k-m+\ell}$, proving \eqref{mburn:eq:M}.
At a different zero $\sigma$ the denominator is nonzero and
the completed transform has a zero of order $m_\sigma$,
so all indicated derivatives vanish. Multiplying the proposed
inverse on the left gives
\[
 \sum_{k=0}^{m-1}c_{\rho,m-\ell-k}d_{\rho,k-m+j}
       =\delta_{\ell j}.
\]
Indeed the nonzero terms have indices summing to $j-\ell$;
if $j<\ell$ there are none, and otherwise their convolution
is the $(j-\ell)$th coefficient of $d_\rho/d_\rho=1$.
Reversing the columns of $M^\rho$ makes it lower triangular
with diagonal $k!d_{\rho,0}$. The reversal has sign
$(-1)^{m(m-1)/2}$, proving the determinant formula.
Substitution into the bilinear pairings proves both dual
identities. Finally $[f,Y]=\langle\overline Y,f\rangle$
in the stated convention, which proves the last assertion.
\end{proof}

For example the full multiplicity-three block, without suppression
of its second-derivative factorial, is
\[
 M^\rho=\begin{pmatrix}0&0&d_0\\0&d_0&d_1\\
                         2d_0&2d_1&2d_2\end{pmatrix},\qquad
 (M^\rho)^{-1}=\begin{pmatrix}c_2&c_1&c_0/2\\
                           c_1&c_0&0\\c_0&0&0\end{pmatrix}.
\]

Here is the exact conversion to original physical jets. Set
\begin{equation}\label{mburn:eq:pq}
 \begin{gathered}
 p_{\rho,0}=\rho(\rho-1),\quad p_{\rho,1}=2\rho-1,
 \quad p_{\rho,2}=1,\quad p_{\rho,n}=0\ (n>2),\\
 \quad q_{\rho,n}=(-1)^n\left((\rho-1)^{-n-1}-\rho^{-n-1}\right).
 \end{gathered}
\end{equation}
These are the Taylor coefficients of $s(s-1)$ and its reciprocal
at $s=\rho$. If $\varphi_k(f)=(b\mburnMR f)^{(k)}(\rho)$ and
$\theta_j(f)=(\mburnML Tf)^{(j)}(\rho)/j!$, then
\begin{equation}\label{mburn:eq:jet-conversion}
 \theta_j(f)=\sum_{k=0}^j p_{\rho,j-k}\frac{\varphi_k(f)}{k!},
 \qquad
 \varphi_k(f)=k!\sum_{j=0}^kq_{\rho,k-j}\theta_j(f).
\end{equation}
These follow by multiplying Taylor series; the reciprocal formula
also follows directly from
$1/(s(s-1))=1/(s-1)-1/s$.
Write $e_{\rho,n}=\sum_{v=0}^{\min(2,n)}p_{\rho,v}d_{\rho,n-v}$
and $v_{\rho,n}=\sum_{j=0}^nq_{\rho,j}c_{\rho,n-j}$.
The finite physical-jet matrix and its inverse are exactly
\begin{equation}\label{mburn:eq:physical-jet}
 N^\rho_{j\ell}=e_{\rho,j-m+\ell},\qquad
 (N^\rho)^{-1}_{\ell j}=v_{\rho,m-\ell-j}.
\end{equation}
The same reciprocal-convolution proof as above verifies both
identities. Its determinant is
$(-1)^{m(m-1)/2}(\rho(\rho-1)d_{\rho,0})^m$.

In particular the local entries of the original unit
\eqref{mburn:eq:unit} and its inverse, including all other roots, are
\begin{equation}\label{mburn:eq:fullunit}
 (U_h)_{\rho,n}=\sum_{i+j+k=n}p_{\rho,i}d_{\rho,j}\alpha_{\rho,k},
 \quad(\varepsilon_h)_{\rho,0}=(U_h)_{\rho,0}^{-1},\quad
 (\varepsilon_h)_{\rho,n}=-(U_h)_{\rho,0}^{-1}
       \sum_{r=1}^n(U_h)_{\rho,r}(\varepsilon_h)_{\rho,n-r}.
\end{equation}
Let $V_h$ be the confluent matrix from original increasing powers
to divided jets, $(V_h)_{(\rho,j),n}=\binom nj\rho^{n-j}$ for
$n\geq j$ and zero otherwise, $0\leq n<d$. This is invertible:
a polynomial of degree $<d$ with every indicated jet zero is
divisible by $h$ and hence is zero. The full coefficient map is
\begin{equation}\label{mburn:eq:coordinate}
 \operatorname{coeff}(r_h(u))=
 V_h^{-1}\operatorname{diag}_{\rho}
       \bigl((\varepsilon_h)_{\rho,j-k}\mathbf1_{j\geq k}\bigr)u.
\end{equation}
Equations \eqref{mburn:eq:partial}, \eqref{mburn:eq:coordinate} compute a
matrix $C_h$ with $f_{r_h(u)/h}=\sum X_{\rho,\ell}(C_hu)_{\rho,\ell}$.
Independently \eqref{mburn:eq:physical-jet} gives
\begin{equation}\label{mburn:eq:Ch}
 C_h=\operatorname{diag}_{\rho}(N^\rho)^{-1}.
\end{equation}
Both descriptions agree because each prescribes the unique
physical jets of $T f_{r/h}$. This proves, rather than omits,
the cancellation of the other-root coefficients between
\eqref{mburn:eq:partial} and \eqref{mburn:eq:coordinate}; the original
physical unit remains fully present in the original power coordinates.

The coefficients in Burnol's residue formula also follow exactly.
For $G=\mburnMR f$ holomorphic at $\rho$, let
\begin{equation}\label{mburn:eq:residue-coeff}
 \lambda_{\rho,\ell}(f)=
 \sum_{k=0}^{m-\ell}\frac{c_{\rho,m-\ell-k}}{k!}
                 (bG)^{(k)}(\rho)
 = [t^{m-\ell}]\frac{G(\rho+t)}{z_\rho(t)}.
\end{equation}
Then for a parameter $w\ne\rho$,
\begin{equation}\label{mburn:eq:residue}
 \operatorname*{Res}_{s=\rho}
   \frac{G(s)}{\zeta(s)}\frac{\zeta(w)}{w-s}
   =\sum_{\ell=1}^m\lambda_{\rho,\ell}(f)
                                    \frac{\zeta(w)}{(w-\rho)^\ell}.
\end{equation}
To prove it, expand $G/\zeta=\sum_{\ell=1}^m
\lambda_{\rho,\ell}(s-\rho)^{-\ell}+O(1)$ and
$1/(w-s)=\sum_{n\geq0}(s-\rho)^n/(w-\rho)^{n+1}$;
the coefficient of $(s-\rho)^{-1}$ is the displayed sum.
No simple-zero convention or missing sign is used.

\section{The full finite Gram and the unchanged relation minimum}

Fix $N\geq d-1$, put $M=N-d$, and interpret $\mburnPP_{-1}=0$.
Use all original divided-jet coordinates in $E_h$ and original
increasing-power coefficients in $\mburnPP_M$. Let $X_h$ have columns
$X_{\rho,\ell}$, and let $F_M$ have columns
$f_j=D^j\Theta\phi_*$ for $0\leq j\leq M$. The exact
prequotient coordinate maps are
\begin{equation}\label{mburn:eq:LN}
 \begin{aligned}
 L_N:E_h\oplus\mburnPP_M&\longrightarrow\mburnPP_N,&
 L_N(u,Q)&=r_h(u)+hQ,\\
 L_N^{-1}P&=\left(j_h((g/h)P),\ (P-r_h(j_h((g/h)P)))/h\right),\\
 \mathcal R_N&=[T X_hC_h,F_M]=[R_h,F_M],&
 \mathcal R_NL_N^{-1}P&=P(D)F_h.
 \end{aligned}
\end{equation}
Every equality follows from \eqref{mburn:eq:all-degree} and
\eqref{mburn:eq:Rh}; both compositions of $L_N$ and its stated
inverse are the identity. Thus the relation summand remains
exactly $h\mburnPP_{N-d}$, not a Burnol orthogonal complement.

Write $s=1/2+it$ and
$x_{\rho,\ell}(s)=\zeta(s)/(s-\rho)^\ell$, with removable
values at critical-line zeros. The following matrices are literal
convergent integrals on the whole real line:
\begin{equation}\label{mburn:eq:grams}
 \begin{aligned}
 (G_B)_{\alpha\beta}
   &=\int_\mburnR\overline{x_\alpha(s)}x_\beta(s)\,\frac{dt}{2\pi},\\
 (G_T)_{\alpha\beta}
   &=\int_\mburnR|a(s)|^2\overline{x_\alpha(s)}x_\beta(s)\,\frac{dt}{2\pi},\\
 K_{\alpha j}
   &=\int_\mburnR|a(s)|^2\overline{x_\alpha(s)}\zeta(s)s^j\,\frac{dt}{2\pi},\\
 S_{ij}
   &=\int_\mburnR|g(s)|^2\overline{s}^{\,i}s^j\,\frac{dt}{2\pi}.
 \end{aligned}
\end{equation}
Here $\alpha,\beta$ run over every $(\rho,\ell)$; $i,j$ run
over $0,\ldots,M$. By Parseval, $G_B=X_h^*X_h$,
$G_T=X_h^*T^*TX_h$, $K=(TX_h)^*F_M$, and $S=F_M^*F_M$.
The first is Burnol's finite metric; the second is the exact
transferred metric. In particular the factor $|a(s)|^2$ is not
replaced by one.

\begin{mburntheorem}[The original attained metric in transferred coordinates]
\label{mburn:thm:gram}
The full Gram and the quotient minimum are exactly
\begin{equation}\label{mburn:eq:schur}
 \mathbb G_N=\mathcal R_N^*\mathcal R_N
   =\begin{pmatrix}A&C\\mburnC^*&S\end{pmatrix}>0,
 \quad A=C_h^*G_TC_h,\quad C=C_h^*K,\quad
 H_N=A-CS^{-1}C^*>0.
\end{equation}
For $M=-1$, both relation blocks are absent and $H_N=A$.
For $M\geq0$ the unique attained lift and its exact minimum are
\begin{equation}\label{mburn:eq:minimizer}
 Q_*=-S^{-1}C^*u,\qquad
 u^*H_Nu=\min_{Q\in\mburnPP_M}
          \mburnnorm{T X_hC_hu+\Theta Q(D)\phi_*}_2^2.
\end{equation}
Equivalently, for every $P\in\mburnPP_N$,
\begin{equation}\label{mburn:eq:originalmetric}
 \mburnnorm{P(D)F_h}_2^2
  =\int_\mburnR|P(1/2+it)|^2\left|(g/h)(1/2+it)\right|^2
                                                   \frac{dt}{2\pi},
\end{equation}
and \eqref{mburn:eq:minimizer} minimizes this exact form on the
unchanged affine fibre $j_h((g/h)P)=u$.
\end{mburntheorem}
\begin{proof}
All columns lie in $\mburnB$, so their products are integrable
in $dx$, and their Mellin transforms give every integral in
\eqref{mburn:eq:grams} by \eqref{mburn:eq:parseval}. This also proves
\eqref{mburn:eq:originalmetric}, without a tail truncation.
If $\mathcal R_N(u,Q)=0$, its full jet is $u$, hence $u=0$.
Then its transform is $gQ=0$, so $Q=0$ because $g$ is not
identically zero. Thus the full finite Gram is positive
definite and so is its relation block $S$ when present.
Completing the exact quadratic form gives
\[
 (u,Q)^*\mathbb G_N(u,Q)
 =u^*(A-CS^{-1}C^*)u
       +(Q+S^{-1}C^*u)^*S(Q+S^{-1}C^*u).
\]
This proves uniqueness, the minimum, and positivity of the
Schur complement. The maps \eqref{mburn:eq:LN} identify its affine
fibre with the original one, proving the final assertion.
\end{proof}

For completeness the transfer is specified for degrees
$0\leq N<d-1$ too. Euclidean division has $Q=0$ and $R=P$;
let $E_N:\mburnC^{N+1}\to\mburnC^d$ send the power coefficients of
$P$ to the complete partial fractions \eqref{mburn:eq:partial}.
Then the physical source is $TX_hE_N$, its full Gram is
$E_N^*G_TE_N>0$, and its attainable jet subspace is exactly
$\{j_h((g/h)P):P\in\mburnPP_N\}$, with no asserted surjection
onto all $E_h$. This proves all-degree coverage without
silently enlarging a low-degree admissible space.

The exact finite comparison of the two different metrics is
\begin{equation}\label{mburn:eq:metriccomp}
 \begin{gathered}
 0<G_T\preceq\mburnnorm T^2G_B,
 \qquad \lambda_hG_B\preceq G_T\preceq\Lambda_hG_B,\\
 \quad \lambda_h=\lambda_{\min}(G_B^{-1/2}G_TG_B^{-1/2})>0,
 \\
 \Lambda_h=\lambda_{\max}(G_B^{-1/2}G_TG_B^{-1/2}).
 \end{gathered}
\end{equation}
These are finite explicit integral matrices, not universal
constants. Positivity follows from linear independence and
injectivity of $T$; the upper bound follows from its norm.
There is no positive lower constant valid for all such finite
divisors: otherwise it would hold on their union, the span
of Burnol's complete $X$ system, hence by density on all
$L_1$, contradicting Proposition~\ref{mburn:prop:quotient}.
This proves the exact obstruction to a uniform isometry or
uniformly bounded inverse, while retaining the positive
finite comparison \eqref{mburn:eq:metriccomp}.

The fixed-$h$ collapse of the original attained metric is
already established in \cite{mburn:OriginalTheta}, equations
RC7--RC8, and is not a new consequence asserted from Burnol's
minimality. Its exact receiver here is
\[
 H_N=H_Z^{(N-d)},\qquad
 H_{N+1}\preceq H_N,\qquad
 \lim_{N\to\infty}\mburnnorm{H_N}=0.
\]
Indeed RC7 uses the same reference columns $R_h$ and precisely
the relation columns $f_0,\ldots,f_{N-d}$, so its displayed
Schur complement is \eqref{mburn:eq:schur} entry for entry.
RC8 proves the limit from the established density of the
original derivative family in $L^2(dx)$. Thus both the
strictly positive finite attained forms and their known
fixed-divisor limit are retained by the full-factor transfer.

\section{Original additive tensor sources and any fixed admissible map}

The preceding equality is preserved, not merely estimated, by
the original tensor construction. For $k\geq1$ and $N\geq d-1$
let $\mathcal S_{k,N}:\mburnC[S]_{\leq N}\to\mburnPP_N^{\otimes k}$
be the literal polynomial substitution
\begin{equation}\label{mburn:eq:multinomial}
 \mathcal S_{k,N}(S^n)=
 \sum_{n_1+\cdots+n_k=n}\frac{n!}{n_1!\cdots n_k!}
                                  s_1^{n_1}\cdots s_k^{n_k}.
\end{equation}
Put $\mathcal C_{k,N}=(L_N^{-1})^{\otimes k}\mathcal S_{k,N}$
and $\mathcal A_{k,N}=\mathcal R_N^{\otimes k}\mathcal C_{k,N}$.
All these are finite column maps into the original physical
space $L^2((0,\infty)^k,dx_1\cdots dx_k)$. By multiplying
the exact Mellin identities in \eqref{mburn:eq:LN},
\begin{equation}\label{mburn:eq:tensorsource}
 \mburnML^{\otimes k}(\mathcal A_{k,N}P)(s_1,\ldots,s_k)
           =P(s_1+\cdots+s_k)\prod_{i=1}^k(g/h)(s_i).
\end{equation}
Thus its complete original power-coefficient Gram is
\begin{equation}\label{mburn:eq:tensorgram}
 \mathcal H_{k,N}=\mathcal C_{k,N}^*
                    \mathbb G_N^{\otimes k}\mathcal C_{k,N}>0.
\end{equation}
In particular, if $w_h(t)=|(g/h)(1/2+it)|^2/(2\pi)$,
tensor Parseval and Tonelli give exactly
\begin{equation}\label{mburn:eq:convolutionmetric}
 \mburnnorm{\mathcal A_{k,N}P}_2^2
  =\int_{\mburnR^k}|P(k/2+i\textstyle\sum t_i)|^2
                       \prod_iw_h(t_i)\,dt_i
  =\int_\mburnR|P(k/2+it)|^2w_h^{*k}(t)\,dt.
\end{equation}
All moments converge since the one-factor functions lie in
$\mburnB$. Positivity follows because $w_h$ is positive almost
everywhere and a nonzero polynomial cannot vanish on a
whole vertical line.

Let $\chi(S)$ be any already specified original full primary
polynomial of degree $q$, and use its exact remainder map
$J_{\chi,N}:\mburnC[S]_{\leq N}\to\mburnC[S]/\chi$ for $N\geq q-1$.
It remains this map, with its original power coordinates and
kernel $\chi\mburnC[S]_{\leq N-q}$. The actual attained metric is
\begin{equation}\label{mburn:eq:quotientmetric}
 G_{k,N}=\left(J_{\chi,N}\mathcal H_{k,N}^{-1}
                                           J_{\chi,N}^*\right)^{-1}.
\end{equation}
Indeed for any positive matrix $H$ and surjection $J$, the
vector $p_*=H^{-1}J^*(JH^{-1}J^*)^{-1}u$ has $Jp_*=u$.
Every other vector in that fibre is $p_*+z$ with $Jz=0$,
and $z^*Hp_*=0$. Expanding its squared norm proves the
minimum $u^*(JH^{-1}J^*)^{-1}u$ and uniqueness.
Apply this to \eqref{mburn:eq:tensorgram} to prove
\eqref{mburn:eq:quotientmetric} with every original relation.
No claim that the roots of $\chi$ are zeros of $g$ is needed
or made; $\chi$ is the fixed target relation, not the
one-factor divisor $h$.

More generally, for any fixed original injective coefficient
map $I_0$ and surjective coefficient map $\Lambda$, restrictions
and further attained quotients are respectively
$I_0^*G_{k,N}I_0$ and
$(\Lambda G_{k,N}^{-1}\Lambda^*)^{-1}$.
The same preceding proof gives the latter, and substitution
gives the former. Consequently every specified finite
sequence of such maps, including a restricted admissible
polynomial subspace, uses exactly the same source and
relation metric after this Burnol transfer. This is an exact
coordinate identity with the factors \eqref{mburn:eq:grams}, not
a replacement of the source by Burnol's Hilbert metric.

\section{Complete-system consequences in their precise topology}

Theorem~\ref{mburn:thm:T} and Burnol's completeness of the $X$'s imply
\begin{equation}\label{mburn:eq:density}
 T(L_1)\subseteq
 \overline{\operatorname{span}\{T X_{\rho,\ell}\}}^{\mburnB}.
\end{equation}
To prove this, approximate any $f\in L_1$ in its Hilbert
norm by finite $X$ combinations and apply every estimate
\eqref{mburn:eq:strongbound}. Each finite combination is the exact
original canonical section of its physical finite jets by
\eqref{mburn:eq:Rh} and \eqref{mburn:eq:physical-jet}. Thus the same
approximations pass through $q$ to the original strong quotient.
Equation \eqref{mburn:eq:density} does not assert that $T(L_1)$
is closed, or that its strong closure or Hilbert closure is
all of $\mburnB$ or all of $\mburnHH$.

For $f$ in Burnol's dense subspace $\mathcal L_1$ of functions
whose right Mellin transforms decrease faster than every
inverse power on vertical strips, his Theorem \texttt{main}
gives the residue series, grouped between his heights $T_n$,
converging to $f$ in $L_1$ norm. The multiple-zero contribution
is exactly \eqref{mburn:eq:residue-coeff}. Applying the continuous
map $T:L_1\to\mburnB$ proves the strong original-source identity
\begin{equation}\label{mburn:eq:strong-residue}
 Tf=\lim_{n\to\infty}
 \sum_{|\Im\rho|<T_n}\sum_{\ell=1}^{m_\rho}
                         \lambda_{\rho,\ell}(f)T X_{\rho,\ell}
                    \quad\text{in every }p_{b,j}.
\end{equation}
The grouping and the dense-subspace domain are part of this
assertion; no ungrouped convergence or arbitrary-$L_1$
residue expansion is inferred. On each fixed zero block its
coefficients reproduce every physical jet by
\eqref{mburn:eq:jet-conversion}. Completeness and minimality
therefore have a calculated receiver in the original
theta construction, with its original full arithmetic
unit and every finite metric. They do not decide the
real parts of the zeros or evaluate an uncomputed
arithmetic sign or asymptotic allocation.

\begin{thebibliography}{9}
\bibitem{mburn:Burnol} Jean-Fran\c{c}ois Burnol,
\emph{Two complete and minimal systems associated with the zeros
of the Riemann zeta function}, Journal de Th\'eorie des Nombres
de Bordeaux \textbf{16} (2004), 65--94.
Original author source \href{https://arxiv.org/abs/math/0203120v7}
{arXiv:math/0203120v7}, Sections 2--6, particularly
Propositions \texttt{thm:mellinLa}, \texttt{zetaprop},
Theorems \texttt{thmA}, \texttt{thmB}, \texttt{main}, and
Corollary \texttt{cor:main}. The local source read for this
derivation is the original \texttt{main.tex}; its SHA-256 is
\texttt{21cde0aa14cfb593aa294daa42d4ed5b9785c30395a5b9f3c55b47c84aa50339}.
\bibitem{mburn:OriginalTheta} Split-Zero research programme,
\emph{Exact receivers from earlier Split Zero mathematics},
19 September 2026, original-source equations T1--T2, U1--U2,
MI1--MI9, RC1--RC10. The defining objects and all receiving
identities used in this note are reproduced and proved above.
\bibitem{mburn:OriginalMetric} Split-Zero research programme,
\emph{The circle heat source in the original arithmetic quotient
metrics}, 19 September 2026, HM1--HM7, HM15--HM21.
The unchanged all-degree source, affine relation space, and
tensor metric are explicitly reconstructed in
\eqref{mburn:eq:LN}--\eqref{mburn:eq:quotientmetric}.
\bibitem{mburn:ClosedTheta} \emph{The closed range of the original full
theta source with every weight, derivative, and factor retained},
20 September 2026, \texttt{ORIGINAL\_THETA\_CLOSED\_RANGE.tex},
Theorem \texttt{thm:closed}; its human closed-range source is
Ralf Meyer, \emph{A spectral interpretation for the zeros of the
Riemann zeta function}, \href{https://arxiv.org/abs/math/0412277v3}
{arXiv:math/0412277v3}, Theorem 3.3. No use of Hilbert closedness
is made in the present transfer.
\end{thebibliography}

\endgroup

\clearpage\begingroup
\part*{Original-kernel projection and full complex-phase control}
\newcommand{\mprojC}{\mathbb C}
\newcommand{\mprojim}{\operatorname{im}}
\newcommand{\mprojdiag}{\operatorname{diag}}
\newcommand{\mprojrank}{\operatorname{rank}}



\section{The available native bracket and the unchanged original class}
On the original domain $k\equiv1\pmod4$, $k\ge77$, $q=(k+1)^2$
and $N\ge q-1$, retain the original finite packet space $E_k$, its actual native attained
metric $G=G_N$, the fixed original subspace $K=\ker\Lambda$, its original
multiplication $A=M_S$, and the specified original class $x\in E_k$.
MS1--3 \cite{mproj:MS}, proved from the complete original native truncation
TR1--5 \cite{mproj:TR}, give
\[
 G_0=\underline G\preceq G\preceq\kappa G_0,\qquad
 \kappa=(1-\epsilon)^{-1},\quad
 \epsilon=L_N^\sharp/(2J)\in(0,1),\quad J>L_N^\sharp/2.          \tag{PD1}
\]
All entries of $G_0$ are the retained native truncated integrals.
Thus PD1 is available at the displayed actual integer cutoff; no
unproved moment estimate is added here.  Inner products are
$\langle u,v\rangle_G=u^*Gv$, conjugate linear in the first argument.
Let $P_0$ and $P$ be the orthogonal projections onto this same $K$
for $G_0$ and $G$, respectively.  For the unchanged original class set
\[
 a=P_0x,\quad b=(I-P_0)x,\quad
 u=\|a\|_{G_0},\quad v=\|b\|_{G_0},\quad
 e=(P-P_0)x.
                                                                    \tag{PD2}
\]
The two mixed pairings to be controlled, including their phases, are
\[
 \begin{array}{ll}
 z=\langle Px,A(I-P)x\rangle_G,
   &z_0=\langle a,Ab\rangle_{G_0},\\
 w=\langle(I-P)x,APx\rangle_G,
   &w_0=\langle b,Aa\rangle_{G_0}.
 \end{array}                                                       \tag{PD3}
\]
No real part or modulus is taken before the exact differences are
formed.  In particular the original complex class is not replaced
by a singular vector or an independently chosen component.

\section{The sharp projection bound}
Use the explicit coordinate isometry $G_0^{1/2}$, then an orthonormal
basis adapted to $K\oplus K^{\perp_{G_0}}$.  In these coordinates
$G_0=I$ and the actual other metric and projections are
\[
 E=\begin{pmatrix}H&X\\X^*&D\end{pmatrix},\quad I\preceq E\preceq\kappa I,
 \qquad P_0=\begin{pmatrix}I&0\\0&0\end{pmatrix},\quad
 P=\begin{pmatrix}I&T\\0&0\end{pmatrix},\quad T=H^{-1}X.          \tag{PD4}
\]
Indeed the residual of a candidate projection $(a+Tb,0)$ is
$(-Tb,b)$, and its pairing with every $(h,0)\in K$ vanishes exactly
when $HT=X$.  The metric $H$ is positive definite, proving the
stated formula and uniqueness.

For each eigenvalue $t\in[1,\kappa]$ one has
$t^2\le(\kappa+1)t-\kappa$.  Diagonalization of the Hermitian
$E$ therefore gives
$E^2\preceq(\kappa+1)E-\kappa I$.  Its upper-left principal block is
\[
 XX^*\preceq(\kappa+1)H-\kappa I-H^2.
\]
Congruence by $H^{-1}$ yields
\[
 TT^*\preceq(\kappa+1)H^{-1}-\kappa H^{-2}-I
           \preceq\frac{(\kappa-1)^2}{4\kappa}I.                 \tag{PD5}
\]
For the last inequality, the scalar function
$(\kappa+1)/t-\kappa/t^2-1$ on $[1,\kappa]$ is maximized at
$t=2\kappa/(\kappa+1)$, as direct differentiation proves; its
value there is $(\kappa-1)^2/(4\kappa)$.  The resulting exact bound is
\[
 \boxed{\|P-P_0\|_{G_0}\le s,
 \qquad s=\frac{\kappa-1}{2\sqrt\kappa}
          =\frac{\epsilon}{2\sqrt{1-\epsilon}},\qquad
 \|e\|_{G_0}\le s v.}                                           \tag{PD6}
\]
The sharper vector statement follows since $(P-P_0)a=0$ and
$e=Tb\in K$.  Thus
\[
 Px=a+e,\quad(I-P)x=b-e,\quad
 \|a+e\|_{G_0}\le u+sv,\quad
 \|b-e\|_{G_0}^2=v^2+\|e\|_{G_0}^2\le(1+s^2)v^2.              \tag{PD7}
\]
The last equality uses the actual $G_0$-orthogonality of $b$ and $e$.

The constant in PD6 is optimal whenever both subspaces have positive
dimension.  On a two-dimensional subspace, take
\[
 H=\frac{2\kappa}{\kappa+1},\qquad
 X=\frac{(\kappa-1)\sqrt\kappa}{\kappa+1},\qquad
 D=\frac{\kappa^2+1}{\kappa+1}.                                  \tag{PD8}
\]
The resulting positive matrix has trace $\kappa+1$ and determinant
$\kappa$, hence eigenvalues $1,\kappa$; its projection coefficient
$H^{-1}X$ is exactly $s$.  Adding identity blocks handles larger
spaces.  If $K$ or its complement is zero, the two projections
are identical and the bound is zero.  The case $\kappa=1$ also
has $P=P_0$.  Interchanging the metrics and applying the same proof
to $G\preceq\kappa G_0\preceq\kappa G$ additionally gives
$\|P-P_0\|_G\le s$, since scaling a metric does not change its
orthogonal projection.

\section{A proved disk for each original complex mixed pairing}
Write $M=\|A\|_{G_0}$.  In the isometric coordinates of PD4 the
exact expansion of the first difference is
\[
 \begin{split}
 z-z_0={}&\langle e,Ab\rangle_0-\langle a,Ae\rangle_0
              -\langle e,Ae\rangle_0\\
          &+\langle a+e,(E-I)A(b-e)\rangle_0.
 \end{split}                                                      \tag{PD9}
\]
Here $\langle\ ,\ \rangle_0$ is the Euclidean realization of
the original $G_0$ inner product.  The other exact difference is
\[
 \begin{split}
 w-w_0={}&\langle b,Ae\rangle_0-\langle e,Aa\rangle_0
              -\langle e,Ae\rangle_0\\
          &+\langle b-e,(E-I)A(a+e)\rangle_0.
 \end{split}                                                      \tag{PD10}
\]
Since $\|E-I\|\le\kappa-1$, Cauchy--Schwarz and PD7 give the
same proved radius for both complete complex differences:
\[
 \begin{split}
 R(M) =M\bigl[&s v^2+suv+s^2v^2\\
                &+(\kappa-1)(u+sv)\sqrt{1+s^2}\,v\bigr],\\
 \boxed{|z-z_0|\le R(M),\qquad |w-w_0|\le R(M).}
 \end{split}                                                      \tag{PD11}
\]
This verifies every term of the proposed common radius.  In
particular it vanishes when $v=0$, because then both original
mixed pairings themselves vanish.

There is a further exact improvement without changing the original
action.  For every complex scalar $\lambda$, orthogonality gives
\[
 \langle Px,\lambda(I-P)x\rangle_G
 =\langle(I-P)x,\lambda Px\rangle_G=0,
\]
and likewise for $P_0,G_0$.  Hence all four quantities in PD3
are unchanged when the right-hand action is written
$A=(A-\lambda I)+\lambda I$, with the scalar summand explicitly
evaluated as zero.  The entire proof PD9--11 therefore also gives
\[
 |z-z_0|,|w-w_0|\le R(\|A-\lambda I\|_{G_0})
                  \quad\hbox{for each specified }\lambda.       \tag{PD12}
\]
In particular the actual original centre $\lambda=k/2$ is available.
This is an exact cancellation in the original pairings, not an
unrecorded coordinate or source normalization.  The constant can
be minimized over any explicitly computed choices of $\lambda$;
no minimizing scalar or asymptotic value is assumed.

\section{A sharper reverse disk from the full Schur block}
Keep the four blocks of the fixed original $A$ in the $G_0$-orthogonal
decomposition:
\[
 A=\begin{pmatrix}A_{11}&A_{12}\\A_{21}&A_{22}\end{pmatrix},\qquad
 S=D-X^*H^{-1}X.                                                 \tag{PD13}
\]
The inverse-block identity gives $S^{-1}$ as the lower-right block
of $E^{-1}$.  From $\kappa^{-1}I\preceq E^{-1}\preceq I$ it follows
that $I\preceq S\preceq\kappa I$.  Equivalently, the same inequalities
follow by taking the full minimum of the $E$ norm over the first
coordinate with the second held fixed; either proof retains the
entire off-diagonal block $X$.

The exact multiplication
\[
 E(-Tb,b)=(0,Sb),\qquad
 E(a+Tb,0)=(H(a+Tb),X^*(a+Tb))
\]
now evaluates both original pairings as
\[
 \begin{split}
 w&=b^*S A_{21}(a+Tb),\quad w_0=b^*A_{21}a,\\
 z&=(a+Tb)^*H
       (A_{12}-A_{11}T+TA_{22}-TA_{21}T)b,
       \quad z_0=a^*A_{12}b.
 \end{split}                                                      \tag{PD14}
\]
Thus if $M_{21}=\|A_{21}\|$ in these original $G_0$ coordinates,
\[
 \boxed{|w-w_0|\le
 R_w^{\rm Schur}:=M_{21}\bigl[(\kappa-1)uv+\kappa s v^2\bigr].}
                                                                    \tag{PD15}
\]
Indeed its two exact difference terms are
$b^*(S-I)A_{21}a$ and $b^*SA_{21}Tb$, and their operator bounds
are the two displayed summands.  This improves on use of the
norm of the whole action whenever the actual block is smaller.
The original isometries from the kernel frame and the minimum
observation lift identify $A_{21}$ precisely with MS4's fixed
lower-metric leakage presentation
$\underline Q^{1/2}(\Lambda AI)\underline H_K^{-1/2}$.
Consequently its norm is the actual largest singular value of
that already defined native matrix, not a new auxiliary action.

An optional blockwise radius for the other direction follows from
the second identity in PD14.  For any specified $\lambda\in\mprojC$ put
$M_{11,\lambda}=\|A_{11}-\lambda I\|$,
$M_{22,\lambda}=\|A_{22}-\lambda I\|$ and
$M_{12}=\|A_{12}\|$.  Since the two scalar terms cancel in
$-A_{11}T+TA_{22}$, expansion gives
\[
 \begin{split}
 |z-z_0|\le R_z^{\rm block}:={}&
 M_{12}\bigl[(\kappa-1)uv+\kappa s v^2\bigr]\\
 &+\kappa(u+sv)v
       \bigl[s(M_{11,\lambda}+M_{22,\lambda})+s^2M_{21}\bigr].
 \end{split}                                                      \tag{PD16}
\]
The first line bounds the two terms from changing $a$ and $H$
in $(a+Tb)^*HA_{12}b$; the second line bounds the three remaining
blocks in PD14.  Every term of the fixed action is therefore
either retained in a bound or cancelled by an exact displayed
identity.  The minimum of the applicable proved radii may be used:
\[
 r_z=\min\{R(\|A-\lambda I\|_{G_0}),R_z^{\rm block}\},\qquad
 r_w=\min\{R(\|A-\lambda I\|_{G_0}),R_w^{\rm Schur}\}.          \tag{PD17}
\]
Using only PD11 instead remains valid; none of these comparisons
asserts that a numerically uncomputed radius is small.

\section{Signed combinations and the exact complete-action cancellation}
The pair of proved disks gives, for every specified complex
$\alpha,\beta$,
\[
 \left|\alpha z+\beta\overline w
         -(\alpha z_0+\beta\overline{w_0})\right|
                        \le|\alpha|r_z+|\beta|r_w.              \tag{PD18}
\]
It also bounds the actual complex product, including its sign-sensitive
real and imaginary parts:
\[
 \boxed{|\overline z w-\overline{z_0}w_0|
       \le |w_0|r_z+|z_0|r_w+r_zr_w.}                           \tag{PD19}
\]
This follows by expanding
$\overline{(z_0+\Delta z)}(w_0+\Delta w)$ and retaining all three
error terms.  Taking real or imaginary parts in this inequality
is legitimate after the complex error has been bounded.  For
example, if its right-hand side is denoted by $R_\times$, the
computed real product lies in
$[\Re(\overline{z_0}w_0)-R_\times,
\Re(\overline{z_0}w_0)+R_\times]$.
This is a rigorous enclosure from the actual native bracket; no
claim that this interval has one sign is made without evaluating
its centre and radius.  A subsequent metric-dependent denominator
must retain its own proved enclosure.

Let $A^{\dagger_G}=G^{-1}A^*G$.  Conjugate symmetry gives the exact
complete-action identity
\[
 \boxed{z+\overline w
       =\langle Px,(A+A^{\dagger_G})(I-P)x\rangle_G.}            \tag{PD20}
\]
For the original action retain its actual centre $c=k/2$ and
$\mathcal W_G=A+A^{\dagger_G}-2cI$.  The scalar contribution in
PD20 is zero by the original orthogonality, so
\[
 z+\overline w=\langle Px,\mathcal W_G(I-P)x\rangle_G.           \tag{PD21}
\]
The original rank-two Hermitian-weight identity underlying DO17 \cite{mproj:DO}
gives $\mprojrank\mathcal W_G\le2$.  If its actual orthonormal spectral
vectors and real eigenvalues are $(v_j,\lambda_j)$, $1\le j\le r_G$
with $r_G\le2$, then PD21 is exactly
\[
 z+\overline w
   =\sum_{j=1}^{r_G}\lambda_j
       \langle Px,v_j\rangle_G\langle v_j,(I-P)x\rangle_G.        \tag{PD22}
\]
This is the full original signed cancellation with all class phases
present.  It neither sets either mixed direction to zero nor
substitutes a rank count for these actual pairings.  Its immediate
finite bound is
\[
 |z+\overline w|
 \le\|\mathcal W_G\|_G\|Px\|_G\|(I-P)x\|_G
 \le\tfrac12\|\mathcal W_G\|_G\|x\|_G^2,                      \tag{PD23}
\]
because the two components are $G$-orthogonal and
$2ab\le a^2+b^2$.  The actual rank-two norm must still be
calculated or bounded; its rank alone supplies no numerical size.

\section{Proof coverage}
PD5--8 proves the sharp metric-dependent projection comparison.
PD9--17 evaluates finite disks for the actual mixed pairings,
including a smaller reverse radius from the original Schur block.
PD18--23 gives the exact receivers into their complex product and
complete-action Hermitian cancellation.  All arguments are
finite-dimensional proofs from the actual native bracket MS1--3
and TR1--5.  They introduce no substituted source, period, class,
projection kernel or unresolved integral evaluation.
\begin{thebibliography}{9}
\bibitem{mproj:MS} Split-Zero programme,
\emph{Native truncation bounds for the original mixed-action singular values},
MS1--8, complete source \texttt{NATIVE\_METRIC\_SINGULAR\_VALUE\_BRACKETS.tex}.
The original two quotient minima in MS2--3 are retained in PD1.
\url{https://github.com/KokunoYumeto/zeta-function-research-reader/tree/405dd5b5d3ab99c5622e1aedb5a16404f6bf94eb/workbenches/splitzero-tandem/continuations/20260920-original-observation-source-kernel}.
\bibitem{mproj:TR} Split-Zero programme,
\emph{Native resolvent truncation at growing degree and the original conductor minimum},
TR1--10, complete source \texttt{NATIVE\_TRUNCATION\_CONDUCTOR\_RETURN.tex}.
\url{https://github.com/KokunoYumeto/zeta-function-research-reader/tree/ab45e221579a0d48ce2885b4ecdf1a6aefc24a20/workbenches/splitzero-tandem/continuations/20260920-native-parity-conductor-return}.
The human multiplier source in its companion is A.~I.~Aptekarev,
G.~L\'opez Lagomasino and A.~Mart\'inez--Finkelshtein,
\emph{Strong asymptotics for the Pollaczek multiple orthogonal polynomials ensembles},
arXiv:1410.1261v1 (2014). The full programme TR proof was read;
this receiver does not claim another reading of that full author paper.
\bibitem{mproj:DO} Split-Zero programme,
\emph{Constant-depth original observation},
DO12--17, complete source \texttt{CONSTANT\_DEPTH\_ORIGINAL\_OBSERVATION.tex},
in the same published edition as MS. DO17 displays the complete original
rank-two attained-source identity before taking its mixed block.
\end{thebibliography}

\endgroup

\clearpage\begingroup
\part*{Independent complete projection derivation}
\newcommand{\mprindC}{\mathbb C}
\newcommand{\mprindim}{\operatorname{im}}
\newcommand{\mprinddiag}{\operatorname{diag}}
\newcommand{\mprindrank}{\operatorname{rank}}



\section{The available native bracket and the unchanged original class}
Retain the original finite packet space $E_k$, its actual native attained
metric $G=G_N$, the fixed original subspace $K=\ker\Lambda$, its original
multiplication $A=M_S$, and the specified original class $x\in E_k$.
MS1--3, proved from the complete original native truncation TR1--5, give
\[
 G_0=\underline G\preceq G\preceq\kappa G_0,\qquad
 \kappa=(1-\epsilon)^{-1},\quad
 \epsilon=L_N^\sharp/(2J)\in(0,1),\quad J>L_N^\sharp/2.          \tag{PD1}
\]
All entries of $G_0$ are the retained native truncated integrals.
Thus PD1 is available at the displayed actual integer cutoff; no
unproved moment estimate is added here.  Inner products are
$\langle u,v\rangle_G=u^*Gv$, conjugate linear in the first argument.
Let $P_0$ and $P$ be the orthogonal projections onto this same $K$
for $G_0$ and $G$, respectively.  For the unchanged original class set
\[
 a=P_0x,\quad b=(I-P_0)x,\quad
 u=\|a\|_{G_0},\quad v=\|b\|_{G_0},\quad
 e=(P-P_0)x.
                                                                    \tag{PD2}
\]
The two mixed pairings to be controlled, including their phases, are
\[
 \begin{array}{ll}
 z=\langle Px,A(I-P)x\rangle_G,
   &z_0=\langle a,Ab\rangle_{G_0},\\
 w=\langle(I-P)x,APx\rangle_G,
   &w_0=\langle b,Aa\rangle_{G_0}.
 \end{array}                                                       \tag{PD3}
\]
No real part or modulus is taken before the exact differences are
formed.  In particular the original complex class is not replaced
by a singular vector or an independently chosen component.

\section{The sharp projection bound}
Use the explicit coordinate isometry $G_0^{1/2}$, then an orthonormal
basis adapted to $K\oplus K^{\perp_{G_0}}$.  In these coordinates
$G_0=I$ and the actual other metric and projections are
\[
 E=\begin{pmatrix}H&X\\X^*&D\end{pmatrix},\quad I\preceq E\preceq\kappa I,
 \qquad P_0=\begin{pmatrix}I&0\\0&0\end{pmatrix},\quad
 P=\begin{pmatrix}I&T\\0&0\end{pmatrix},\quad T=H^{-1}X.          \tag{PD4}
\]
Indeed the residual of a candidate projection $(a+Tb,0)$ is
$(-Tb,b)$, and its pairing with every $(h,0)\in K$ vanishes exactly
when $HT=X$.  The metric $H$ is positive definite, proving the
stated formula and uniqueness.

For each eigenvalue $t\in[1,\kappa]$ one has
$t^2\le(\kappa+1)t-\kappa$.  Diagonalization of the Hermitian
$E$ therefore gives
$E^2\preceq(\kappa+1)E-\kappa I$.  Its upper-left principal block is
\[
 XX^*\preceq(\kappa+1)H-\kappa I-H^2.
\]
Congruence by $H^{-1}$ yields
\[
 TT^*\preceq(\kappa+1)H^{-1}-\kappa H^{-2}-I
           \preceq\frac{(\kappa-1)^2}{4\kappa}I.                 \tag{PD5}
\]
For the last inequality, the scalar function
$(\kappa+1)/t-\kappa/t^2-1$ on $[1,\kappa]$ is maximized at
$t=2\kappa/(\kappa+1)$, as direct differentiation proves; its
value there is $(\kappa-1)^2/(4\kappa)$.  The resulting exact bound is
\[
 \boxed{\|P-P_0\|_{G_0}\le s,
 \qquad s=\frac{\kappa-1}{2\sqrt\kappa}
          =\frac{\epsilon}{2\sqrt{1-\epsilon}},\qquad
 \|e\|_{G_0}\le s v.}                                           \tag{PD6}
\]
The sharper vector statement follows since $(P-P_0)a=0$ and
$e=Tb\in K$.  Thus
\[
 Px=a+e,\quad(I-P)x=b-e,\quad
 \|a+e\|_{G_0}\le u+sv,\quad
 \|b-e\|_{G_0}^2=v^2+\|e\|_{G_0}^2\le(1+s^2)v^2.              \tag{PD7}
\]
The last equality uses the actual $G_0$-orthogonality of $b$ and $e$.

The constant in PD6 is optimal whenever both subspaces have positive
dimension.  On a two-dimensional subspace, take
\[
 H=\frac{2\kappa}{\kappa+1},\qquad
 X=\frac{(\kappa-1)\sqrt\kappa}{\kappa+1},\qquad
 D=\frac{\kappa^2+1}{\kappa+1}.                                  \tag{PD8}
\]
The resulting positive matrix has trace $\kappa+1$ and determinant
$\kappa$, hence eigenvalues $1,\kappa$; its projection coefficient
$H^{-1}X$ is exactly $s$.  Adding identity blocks handles larger
spaces.  If $K$ or its complement is zero, the two projections
are identical and the bound is zero.  The case $\kappa=1$ also
has $P=P_0$.  Interchanging the metrics and applying the same proof
to $G\preceq\kappa G_0\preceq\kappa G$ additionally gives
$\|P-P_0\|_G\le s$, since scaling a metric does not change its
orthogonal projection.

\section{A proved disk for each original complex mixed pairing}
Write $M=\|A\|_{G_0}$.  In the isometric coordinates of PD4 the
exact expansion of the first difference is
\[
 \begin{split}
 z-z_0={}&\langle e,Ab\rangle_0-\langle a,Ae\rangle_0
              -\langle e,Ae\rangle_0\\
          &+\langle a+e,(E-I)A(b-e)\rangle_0.
 \end{split}                                                      \tag{PD9}
\]
Here $\langle\ ,\ \rangle_0$ is the Euclidean realization of
the original $G_0$ inner product.  The other exact difference is
\[
 \begin{split}
 w-w_0={}&\langle b,Ae\rangle_0-\langle e,Aa\rangle_0
              -\langle e,Ae\rangle_0\\
          &+\langle b-e,(E-I)A(a+e)\rangle_0.
 \end{split}                                                      \tag{PD10}
\]
Since $\|E-I\|\le\kappa-1$, Cauchy--Schwarz and PD7 give the
same proved radius for both complete complex differences:
\[
 \begin{split}
 R(M) =M\bigl[&s v^2+suv+s^2v^2\\
                &+(\kappa-1)(u+sv)\sqrt{1+s^2}\,v\bigr],\\
 \boxed{|z-z_0|\le R(M),\qquad |w-w_0|\le R(M).}
 \end{split}                                                      \tag{PD11}
\]
This verifies every term of the proposed common radius.  In
particular it vanishes when $v=0$, because then both original
mixed pairings themselves vanish.

There is a further exact improvement without changing the original
action.  For every complex scalar $\lambda$, orthogonality gives
\[
 \langle Px,\lambda(I-P)x\rangle_G
 =\langle(I-P)x,\lambda Px\rangle_G=0,
\]
and likewise for $P_0,G_0$.  Hence all four quantities in PD3
are unchanged when the right-hand action is written
$A=(A-\lambda I)+\lambda I$, with the scalar summand explicitly
evaluated as zero.  The entire proof PD9--11 therefore also gives
\[
 |z-z_0|,|w-w_0|\le R(\|A-\lambda I\|_{G_0})
                  \quad\hbox{for each specified }\lambda.       \tag{PD12}
\]
In particular the actual original centre $\lambda=k/2$ is available.
This is an exact cancellation in the original pairings, not an
unrecorded coordinate or source normalization.  The constant can
be minimized over any explicitly computed choices of $\lambda$;
no minimizing scalar or asymptotic value is assumed.

\section{A sharper reverse disk from the full Schur block}
Keep the four blocks of the fixed original $A$ in the $G_0$-orthogonal
decomposition:
\[
 A=\begin{pmatrix}A_{11}&A_{12}\\A_{21}&A_{22}\end{pmatrix},\qquad
 S=D-X^*H^{-1}X.                                                 \tag{PD13}
\]
The inverse-block identity gives $S^{-1}$ as the lower-right block
of $E^{-1}$.  From $\kappa^{-1}I\preceq E^{-1}\preceq I$ it follows
that $I\preceq S\preceq\kappa I$.  Equivalently, the same inequalities
follow by taking the full minimum of the $E$ norm over the first
coordinate with the second held fixed; either proof retains the
entire off-diagonal block $X$.

The exact multiplication
\[
 E(-Tb,b)=(0,Sb),\qquad
 E(a+Tb,0)=(H(a+Tb),X^*(a+Tb))
\]
now evaluates both original pairings as
\[
 \begin{split}
 w&=b^*S A_{21}(a+Tb),\quad w_0=b^*A_{21}a,\\
 z&=(a+Tb)^*H
       (A_{12}-A_{11}T+TA_{22}-TA_{21}T)b,
       \quad z_0=a^*A_{12}b.
 \end{split}                                                      \tag{PD14}
\]
Thus if $M_{21}=\|A_{21}\|$ in these original $G_0$ coordinates,
\[
 \boxed{|w-w_0|\le
 R_w^{\rm Schur}:=M_{21}\bigl[(\kappa-1)uv+\kappa s v^2\bigr].}
                                                                    \tag{PD15}
\]
Indeed its two exact difference terms are
$b^*(S-I)A_{21}a$ and $b^*SA_{21}Tb$, and their operator bounds
are the two displayed summands.  This improves on use of the
norm of the whole action whenever the actual block is smaller.
The original isometries from the kernel frame and the minimum
observation lift identify $A_{21}$ precisely with MS4's fixed
lower-metric leakage presentation
$\underline Q^{1/2}(\Lambda AI)\underline H_K^{-1/2}$.
Consequently its norm is the actual largest singular value of
that already defined native matrix, not a new auxiliary action.

An optional blockwise radius for the other direction follows from
the second identity in PD14.  For any specified $\lambda\in\mprindC$ put
$M_{11,\lambda}=\|A_{11}-\lambda I\|$,
$M_{22,\lambda}=\|A_{22}-\lambda I\|$ and
$M_{12}=\|A_{12}\|$.  Since the two scalar terms cancel in
$-A_{11}T+TA_{22}$, expansion gives
\[
 \begin{split}
 |z-z_0|\le R_z^{\rm block}:={}&
 M_{12}\bigl[(\kappa-1)uv+\kappa s v^2\bigr]\\
 &+\kappa(u+sv)v
       \bigl[s(M_{11,\lambda}+M_{22,\lambda})+s^2M_{21}\bigr].
 \end{split}                                                      \tag{PD16}
\]
The first line bounds the two terms from changing $a$ and $H$
in $(a+Tb)^*HA_{12}b$; the second line bounds the three remaining
blocks in PD14.  Every term of the fixed action is therefore
either retained in a bound or cancelled by an exact displayed
identity.  The minimum of the applicable proved radii may be used:
\[
 r_z=\min\{R(\|A-\lambda I\|_{G_0}),R_z^{\rm block}\},\qquad
 r_w=\min\{R(\|A-\lambda I\|_{G_0}),R_w^{\rm Schur}\}.          \tag{PD17}
\]
Using only PD11 instead remains valid; none of these comparisons
asserts that a numerically uncomputed radius is small.

\section{Signed combinations and the exact complete-action cancellation}
The pair of proved disks gives, for every specified complex
$\alpha,\beta$,
\[
 \left|\alpha z+\beta\overline w
         -(\alpha z_0+\beta\overline{w_0})\right|
                        \le|\alpha|r_z+|\beta|r_w.              \tag{PD18}
\]
It also bounds the actual complex product, including its sign-sensitive
real and imaginary parts:
\[
 \boxed{|\overline z w-\overline{z_0}w_0|
       \le |w_0|r_z+|z_0|r_w+r_zr_w.}                           \tag{PD19}
\]
This follows by expanding
$\overline{(z_0+\Delta z)}(w_0+\Delta w)$ and retaining all three
error terms.  Taking real or imaginary parts in this inequality
is legitimate after the complex error has been bounded.  For
example, if its right-hand side is denoted by $R_\times$, the
computed real product lies in
$[\Re(\overline{z_0}w_0)-R_\times,
\Re(\overline{z_0}w_0)+R_\times]$.
This is a rigorous enclosure from the actual native bracket; no
claim that this interval has one sign is made without evaluating
its centre and radius.  A subsequent metric-dependent denominator
must retain its own proved enclosure.

Let $A^{\dagger_G}=G^{-1}A^*G$.  Conjugate symmetry gives the exact
complete-action identity
\[
 \boxed{z+\overline w
       =\langle Px,(A+A^{\dagger_G})(I-P)x\rangle_G.}            \tag{PD20}
\]
For the original action retain its actual centre $c=k/2$ and
$\mathcal W_G=A+A^{\dagger_G}-2cI$.  The scalar contribution in
PD20 is zero by the original orthogonality, so
\[
 z+\overline w=\langle Px,\mathcal W_G(I-P)x\rangle_G.           \tag{PD21}
\]
The original rank-two Hermitian-weight identity underlying DO17
gives $\mprindrank\mathcal W_G\le2$.  If its actual orthonormal spectral
vectors and real eigenvalues are $(v_j,\lambda_j)$, $1\le j\le r_G$
with $r_G\le2$, then PD21 is exactly
\[
 z+\overline w
   =\sum_{j=1}^{r_G}\lambda_j
       \langle Px,v_j\rangle_G\langle v_j,(I-P)x\rangle_G.        \tag{PD22}
\]
This is the full original signed cancellation with all class phases
present.  It neither sets either mixed direction to zero nor
substitutes a rank count for these actual pairings.  Its immediate
finite bound is
\[
 |z+\overline w|
 \le\|\mathcal W_G\|_G\|Px\|_G\|(I-P)x\|_G
 \le\tfrac12\|\mathcal W_G\|_G\|x\|_G^2,                      \tag{PD23}
\]
because the two components are $G$-orthogonal and
$2ab\le a^2+b^2$.  The actual rank-two norm must still be
calculated or bounded; its rank alone supplies no numerical size.

\section{Proof coverage}
PD5--8 proves the sharp metric-dependent projection comparison.
PD9--17 evaluates finite disks for the actual mixed pairings,
including a smaller reverse radius from the original Schur block.
PD18--23 gives the exact receivers into their complex product and
complete-action Hermitian cancellation.  All arguments are
finite-dimensional proofs from the actual native bracket MS1--3
and TR1--5.  They introduce no substituted source, period, class,
projection kernel or unresolved integral evaluation.

\endgroup

\clearpage\begingroup
\part*{Complete moment-to-complex-current composition}
\newcommand{\mcompC}{\mathbb C}
\newcommand{\mcomptr}{\operatorname{tr}}


\section{Original data and the explicitly proved moment margin}
For the arithmetic specialization used here, retain the original simple
quartet packet, $k\equiv1\pmod4$, $k\ge77$, $q=(k+1)^2$ and
$N\ge q-1$, exactly as in the complete precision and projection
proofs below. The finite matrix argument applies to their admitted
maps; it does not enlarge the proved arithmetic lower-envelope domain.
This calculation composes the full native precision proof
\cite{mcomp:precision} with the original-metric projection calculation
\cite{mcomp:projection}. It does not select a new source, quotient, class or
projection kernel. Fix the original finite polynomial source degree $N$,
its actual parity-coordinate Gram matrix $H$, its original onto map
$J$ to the finite arithmetic packet $E_k$, the original multiplication
$A=M_S$, the specified original class $x$, and $K=\ker\Lambda\subset E_k$.
The map $J$ includes the retained physical unit and original coordinate
conversion. In particular it is not replaced by a coordinate projection.
The actual quotient metric is
\[
 G=(JH^{-1}J^*)^{-1}.                                      \tag{MC1}
\]
All inner products are conjugate linear in the first variable.

Here is the concrete positive lower margin already proved from the actual
arithmetic measure in \cite{mcomp:precision}. Write its two parity blocks as
moment matrices of degrees $n_0,n_1$, omitting an empty block. If the
retained density lower bound on $[1,2]$ is $d_\nu>0$, put
\[
 (L_t)_{ij}=\frac{2^{i+j+1}-1}{i+j+1}\quad(0\le i,j\le t),
 \qquad g_t=\frac{\det L_t}{(\mcomptr L_t)^t},\qquad
 h=d_\nu\min_{b}g_{n_b},\qquad d=\max_b(n_b+1).              \tag{MC2}
\]
The formula
\[
 \det L_t=\prod_{j=0}^t\frac{(j!)^4}{(2j)!(2j+1)!}
\]
and positivity of the integral Gram on $[1,2]$ give $h>0$.
For the odd parity block the extra factor $u$ on this interval is at
least one, so the same $d_\nu$ is valid. The moment proof gives
$hI\preceq H$. Choose the explicit upper cap $U$ as the sum of
the proved upper caps for every diagonal moment of both blocks. Since
$H$ is positive, $H\preceq(\mcomptr H)I\preceq UI$.
Neither $d_\nu$ nor any total mass is set equal to one.

If each required real moment is enclosed with absolute error at most
$\varepsilon$, a real symmetric midpoint matrix $\widehat H$ has
$\|H-\widehat H\|\le d\varepsilon$: each block has row sum at most
its size times $\varepsilon$. Set $\delta=d\varepsilon$ and
\[
 H_-=\widehat H-\delta I,\quad H_+=\widehat H+\delta I,
 \qquad \theta=\delta/h\le\tfrac14.
                                                               \tag{MC3}
\]
Then
\[
 (1-2\theta)H\preceq H-2\delta I\preceq H_-
 \preceq H\preceq H_+,
 \qquad H_-\succeq\tfrac h2 I.                                \tag{MC4}
\]
In particular all inverses below have a proved positive margin. The
comparison $(1-2\theta)H\preceq H_-$ follows from $hI\preceq H$;
it does not require assuming a positive midpoint.

\section{The unchanged quotient and a computable comparison norm}
Define
\[
 G_0=(JH_-^{-1}J^*)^{-1},\qquad W=(JJ^*)^{-1},\qquad
 \kappa=(1-2\theta)^{-1}.
\]
The map $J$ is onto, so $W$ exists. Order reversal under positive
inversion and congruence by this same $J$ give
\[
 G_0\preceq G\preceq\kappa G_0,\qquad
 \tfrac h2 W\preceq G_0\preceq UW.                            \tag{MC5}
\]
The last upper bound uses $H_-\preceq H\preceq UI$, not an
independently enlarged source. The matrix $W$ is only a comparison
used to evaluate constants; the metric in the current remains $G$.
For an operator $B$ write
$\|B\|_W^2=\sup_{v\ne0}v^*B^*WBv/(v^*Wv)$. This is a finite
generalized eigenvalue, retaining the actual $J$ and $B$.
For the original $A$ set
\[
 \overline M=2\sqrt{U/h}\,\|A\|_W,\qquad
 C=\overline M\,U\,\|x\|_W^2.                                \tag{MC6}
\]
Equation MC5 proves $\|A\|_{G_0}\le\sqrt{2U/h}\|A\|_W
\le\overline M$ and $\|x\|_{G_0}^2\le U\|x\|_W^2$.
Thus all constants in MC6 are evaluated from the original finite maps
and proved source moment caps. They do not depend on a chosen
maximizing class.

\section{Exact projection and complex differences}
Let $P_0,P$ be the $G_0,G$ orthogonal projections onto the same $K$.
Put $a=P_0x$, $b=(I-P_0)x$, $u=\|a\|_{G_0}$, $v=\|b\|_{G_0}$.
The explicit isometry $G_0^{1/2}$ and an orthonormal basis of its
image of $K$ represent the relative metric as
\[
 E=\begin{pmatrix}B&Y\\Y^*&D\end{pmatrix},\quad
 I\preceq E\preceq\kappa I,
 \qquad P=\begin{pmatrix}I&B^{-1}Y\\0&0\end{pmatrix}.
\]
This is a displayed coordinate calculation, not a change of the
physical source. Since
$E^2\preceq(\kappa+1)E-\kappa I$, its upper-left block yields
\[
 B^{-1}YY^*B^{-1}\preceq
 (\kappa+1)B^{-1}-\kappa B^{-2}-I
 \preceq\frac{(\kappa-1)^2}{4\kappa}I.
\]
The last scalar maximum occurs at $2\kappa/(\kappa+1)$.
Consequently $e=(P-P_0)x\in K$ satisfies
\[
 \|e\|_{G_0}\le s v,\quad s=\frac{\theta}{\sqrt{1-2\theta}},
 \quad Px=a+e,\quad (I-P)x=b-e,\quad
 \|b-e\|_{G_0}\le\sqrt{1+s^2}\,v.                            \tag{MC7}
\]
Define the two full complex pairings
\[
 z=\langle Px,A(I-P)x\rangle_G,\quad
 w=\langle(I-P)x,APx\rangle_G,
 \quad z_0=\langle a,Ab\rangle_{G_0},\quad
 w_0=\langle b,Aa\rangle_{G_0}.                               \tag{MC8}
\]
In the isometric coordinates just displayed their exact differences are
\begin{align*}
z-z_0={}&\langle e,Ab\rangle_0-\langle a,Ae\rangle_0
 -\langle e,Ae\rangle_0+\langle a+e,(E-I)A(b-e)\rangle_0,\\
w-w_0={}&\langle b,Ae\rangle_0-\langle e,Aa\rangle_0
 -\langle e,Ae\rangle_0+\langle b-e,(E-I)A(a+e)\rangle_0.
\end{align*}
Writing $M=\|A\|_{G_0}$, Cauchy--Schwarz bounds both absolute
differences by the complete radius
\[
 R=M\{sv^2+suv+s^2v^2+
 (\kappa-1)(u+sv)\sqrt{1+s^2}\,v\}.                          \tag{MC9}
\]
This rederives the cited projection proof's PD11 and keeps the metric
variation term as well as the motion of each vector.

For $0\le\theta\le1/4$, use
$s\le\sqrt2\theta$, $s^2\le2\theta^2$,
$\kappa-1\le4\theta$, and
$\sqrt{1+s^2}\le\sqrt{9/8}=3/(2\sqrt2)$.
Substitution in MC9 gives
\[
 R\le M\theta\{4\sqrt2\,uv+(2+\sqrt2)v^2\}
 \le6M\theta(u^2+v^2)\le6C\theta.                           \tag{MC10}
\]
For clarity, the $v^2$ coefficient after division by $\theta$ is
at most $\sqrt2+2\theta+6\theta\le\sqrt2+2$.
The second inequality follows from the identity
\[
 6(u^2+v^2)-4\sqrt2\,uv-(2+\sqrt2)v^2
 =6\left(u-\frac{\sqrt2}{3}v\right)^2
   +\left(\frac83-\sqrt2\right)v^2\ge0.
\]
Also $|z_0|,|w_0|\le Muv\le M\|x\|_{G_0}^2/2\le C/2$.
No real part has replaced either complex pairing.

\section{A direct moment tolerance for the original product}
Expand the full product, retaining its mixed error:
\[
 \overline z w-\overline{z_0}w_0
 =\overline{(z-z_0)}w_0+\overline{z_0}(w-w_0)
   +\overline{(z-z_0)}(w-w_0).
\]
Equations MC10 and $\theta^2\le\theta/4$ therefore prove
\[
 \boxed{\left|\overline z w-\overline{z_0}w_0\right|
 \le 6C^2\theta+36C^2\theta^2\le15C^2\theta.}               \tag{MC11}
\]
Given a positive absolute tolerance $t_*$ and $C>0$, the explicit
finite moment tolerance
\[
 \boxed{\theta_* =\min\{1/4,t_*/(15C^2)\},\qquad
             \varepsilon\le\theta_*h/d}                    \tag{MC12}
\]
suffices. If $C=0$, either the original class is zero or the original
action has zero operator norm; both pairings then vanish exactly.
For $K=0$ or $K=E_k$ they also vanish exactly. The real and imaginary
parts each inherit MC11 only after the complex difference has been
bounded. No sign is assigned to an interval still containing zero.

The one-factor arithmetic input remains
$\eta_\ell=\int y^\ell w_h(y)\,dy$, with convolution moments
\[
 \zeta_L=\sum_{\ell_1+\cdots+\ell_k=L}
 \frac{L!}{\ell_1!\cdots\ell_k!}\prod_{i=1}^k\eta_{\ell_i}.
\]
If $|\eta_\ell|\le B$ through $L\le2r$ and the one-factor errors
are at most $\tau\le1$, telescoping the difference of each product
gives error at most $k\tau(B+1)^{k-1}$ per product. The sum of the
multinomial coefficients is $k^L$, so it suffices to take
\[
 \tau\le\min\left\{1,
 \frac{\varepsilon}{k^{2r+1}(B+1)^{k-1}}\right\}.             \tag{MC13}
\]
These are the actual mass-bearing moments, not moments of a substitute
measure. The exponential-tail argument and original theta evaluation
in \cite{mcomp:precision} supply finite integral enclosures for them.
Numerical roundoff in inverses, projections and complex arithmetic
must be enclosed in addition to MC11; exact matrix formulas do not
license ignoring that error.

This proves a finite precision prescription for the original complex
current. It assigns no numerical values to the native integrals, no
uniform growing-degree rate, and no conclusion about the Riemann
hypothesis. The original measure and all retained complex phases remain
in every displayed map.
\begin{thebibliography}{9}
\bibitem{mcomp:precision}Programme calculation, \emph{Native moment precision},
\texttt{NATIVE\_MOMENT\_PRECISION.tex},20 September2026,
complete source accompanying this edition, equations for interval
Gram floors, original quotient transport and one-factor moment precision.
Human inputs and their precise uses are cited in that full proof.
\bibitem{mcomp:projection}Programme calculation, \emph{Sharp original-metric
projection control and complex mixed-current return},
\texttt{PROJECTION\_PHASE\_RETURN.tex},20 September2026, PD1--PD23,
complete source accompanying this edition. The sharper constant15 in
MC11 was also obtained by the independent moment--phase derivation
accompanying this edition; it replaces the root's initial valid constant192.
\end{thebibliography}

\endgroup

\clearpage\begingroup
\part*{Independent complete composition with the stronger constant15}
\newcommand{\mcindC}{\mathbb C}
\newcommand{\mcindtr}{\operatorname{tr}}
\newcommand{\mcinddiag}{\operatorname{diag}}
\newcommand{\mcindnorm}[1]{\left\lVert#1\right\rVert}


This independent calculation verifies the proposed constant $192$
and strengthens it to $15$ with exactly the same original constant
$C$ defined below. The complete PD1--23 proof \cite{mcind:PD} was read
through its end. The original native measure, mass, quotient and
observation maps, multiplication, class, kernel and complex phases
are retained. No native scalar value or growing-degree convergence
is asserted.

\section{Actual domain, source constants, and all original maps}

Use precisely the original simple-quartet domain of PD1:
$k\equiv1\pmod4$, $k\ge77$, $q=(k+1)^2$, and $N\ge q-1$.
The original full quartet divisor, completed function, and native
source are
\[
 \rho_*=\tfrac12+\delta+i\gamma,\quad0<\delta<\tfrac12,\quad\gamma>2,
 \qquad
 h(s)=\prod_{z\in\{\rho_*,\bar\rho_*,1-\rho_*,1-\bar\rho_*\}}(s-z),
 \quad g=2\xi,
\]
\[
 w_h(y)=\frac{|(g/h)(1/2+iy)|^2}{2\pi},\quad m_k=w_h^{*k},
 \quad \mu_h=\int_{\mathbb R}w_h(y)dy,\quad c=k/2,\quad S=c+iy.
\]
The native parity measures are
$d\nu_0(x)=m_k(\sqrt x)x^{-1/2}dx$, $d\nu_1=x\,d\nu_0$.
In particular their original masses are $\mu_h^k$ and
$k\eta_2\mu_h^{k-1}$, with $\eta_2=\int y^2w_h(y)dy$.
Put $n_0=\lfloor N/2\rfloor$, $n_1=\lfloor(N-1)/2\rfloor$,
$d_{\rm src}=\max(n_0+1,n_1+1)$, and retain the exact native
parity Gram
\begin{equation}\tag{MC1}\label{mcind:eq:H}
 H=\mcinddiag\left([\int x^{i+j}d\nu_0]_{i,j=0}^{n_0},
               [\int x^{i+j}d\nu_1]_{i,j=0}^{n_1}\right).
\end{equation}

Here are the established explicit original constants from
\cite{mcind:NP,mcind:CAI}, specialized only to the original simple-quartet
multiplicity $m=1$:
\[
\begin{gathered}
 \alpha=\pi/2,\quad c_\Gamma=\sqrt2e^{-7/6},\quad
 C_\Gamma=\sqrt{2\sqrt5}e^{2/3},\quad M=25,\quad B_h=50,\\
 \vartheta_h=\int_{-1}^1w_h(y)dy,\quad
 M_\alpha=\int_{\mathbb R}e^{\alpha y}w_h(y)dy,\\
 a_k=\frac{c_h\vartheta_h^{k-3}e^{-\alpha(k-3)}(k-2)^{-50}}
                 {C_\Gamma2^{25}},\qquad
 A_k=\frac{C_h^{\rm tilt}}{c_\Gamma}k^{9/2}M_\alpha^{k-1},\\
 d_\nu=\frac{a_kc_\Gamma3^{-25}e^{-\alpha\sqrt2}}
                 {\sqrt2\sqrt{1+\sqrt2}},\qquad
 U_{\epsilon,r}=\frac{2A_kC_\Gamma[2(r+\epsilon)]!}
                         {\alpha^{2(r+\epsilon)+1}}.
\end{gathered}
\]
The TW7 and BT12 constants $c_h,C_h^{\rm tilt}$ and the two
displayed original scalar integrals are not assigned new values.
The native density on $[1,2]$ is at least $d_\nu$ in both parity
blocks: this follows from
$m_k(y)\ge a_k(1+y^2)^{-25}
c_\Gamma e^{-\alpha|y|}(1+|y|)^{-1/2}$ and the exact
factor $x^{-1/2}$. The upper moment cap follows by integrating
$2A_kC_\Gamma y^{2(r+\epsilon)}e^{-\alpha y}$ on the positive
half-line, as proved in \cite{mcind:NP}.

For $t\ge0$ set
\[
 J_t^{\rm int}=\left[\frac{2^{i+j+1}-1}{i+j+1}\right]_{i,j=0}^t,
 \quad g_t=\frac{\det J_t^{\rm int}}{(\mcindtr J_t^{\rm int})^t},
 \quad
 \det J_t^{\rm int}=\prod_{j=0}^t
                   \frac{(j!)^4}{(2j)!(2j+1)!}.
\]
These are exact rational constants, with $g_0=1$. The full
determinant identity is proved in \cite{mcind:NP}. Positivity follows
from the integral of a nonzero polynomial's squared modulus on
$[1,2]$; and, if $\lambda_1$ is the least eigenvalue,
$\det J_t^{\rm int}\le\lambda_1(\mcindtr J_t^{\rm int})^t$.
Consequently the following are actual native margins, not new
positivity assumptions:
\begin{equation}\tag{MC2}\label{mcind:eq:margins}
 h_*=d_\nu\min(g_{n_0},g_{n_1}),\qquad
 U=\sum_{\epsilon=0}^1\sum_{i=0}^{n_\epsilon}U_{\epsilon,2i},
 \qquad 0<h_*I\preceq H\preceq UI.
\end{equation}
The scalar $h_*$ is the $h$ in the proposed tolerance; the star
only distinguishes it from the already fixed divisor $h(s)$.

Let $J$ be the unchanged onto quotient map from these parity
coefficients to the original packet space $E_k$. In the original
physical-unit $S$ presentation it is exactly
\begin{equation}\tag{MC3}\label{mcind:eq:Jmap}
 \begin{gathered}
 J=J_{S,N}L_N^{-1},\quad L_N=\Pi_NT_N,\quad
 (T_N)_{rj}=\binom jr c^{j-r}i^r,\quad
 (T_N^{-1})_{rj}=\binom jr(-c)^{j-r}i^{-j}\quad(r\le j),\\
 \det L_N=(\det\Pi_N)i^{N(N+1)/2}.
 \end{gathered}
\end{equation}
Here $\Pi_N$ places even powers before odd powers and the other
triangular entries vanish. Expansion of $(c+iy)^j$ and
$i^{-j}(S-c)^j$ proves the two maps, including the determinant
phase. The original physical unit remains in $J_{S,N}$.
This $J$ is the quotient to $E_k$, not the composite observation
$\Lambda J$. Keep $K=\ker\Lambda\subset E_k$,
$A=M_S:E_k\to E_k$, and the same specified original class
$x\in E_k$. In particular $\ker J$ and $K$ are distinct fibres
and are not merged.
Neither the original action nor the specified original class is
recomputed from the midpoint Gram; in particular no midpoint
orthogonal-polynomial class replaces $x$.

\section{Moment error to the original class metric}

Choose real certified moment midpoints with error at most
$\varepsilon$ in every required original parity entry, using the
literal parity shift for their common base sequence. Their block
Hankel matrix $\widehat H$ satisfies
$\mcindnorm{H-\widehat H}\le d_{\rm src}\varepsilon$, because
each Hermitian block's row sum is at most its dimension times
$\varepsilon$. For $0<\theta\le1/4$ impose
\begin{equation}\tag{MC4}\label{mcind:eq:momentbudget}
 \delta_*=d_{\rm src}\varepsilon\le\theta h_*,\quad
 H_{\rm lo}=\widehat H-\delta_*I,\quad
 H_{\rm hi}=\widehat H+\delta_*I,\quad
 \kappa=(1-2\theta)^{-1}.
\end{equation}
Then
\begin{equation}\tag{MC5}\label{mcind:eq:sourceorder}
 \frac{h_*}{2}I\preceq H_{\rm lo}\preceq H\preceq H_{\rm hi},
 \qquad H_{\rm hi}\preceq\kappa H_{\rm lo},\qquad
 H_{\rm lo}\preceq UI.
\end{equation}
Indeed $H_{\rm lo}\succeq(h_*-2\delta_*)I\succeq h_*I/2$.
The endpoint difference is $2\delta_*I$, so their relative factor
is at most $1+2\delta_* /(h_*-2\delta_*)\le\kappa$.
The last inequality uses the actual order $H_{\rm lo}\preceq
H\preceq UI$; it strengthens the proposed $H_{\rm lo}\preceq2UI$.

Set
\begin{equation}\tag{MC6}\label{mcind:eq:G}
 W=(JJ^*)^{-1},\quad
 G_0=(JH_{\rm lo}^{-1}J^*)^{-1},\quad
 G=(JH^{-1}J^*)^{-1}.
\end{equation}
All inverses exist by \eqref{mcind:eq:sourceorder} and the original
surjectivity of $J$. For every positive source matrix $D$ and
fixed class $y$ the attained minimum is
\[
 \min_{Jb=y}b^*Db=y^*(JD^{-1}J^*)^{-1}y,
 \qquad b_{\min}=D^{-1}J^*(JD^{-1}J^*)^{-1}y.
\]
The displayed vector belongs to the fibre; its $D$ pairing with
every vector of $\ker J$ is zero, which proves the formula by
expanding the full square. Minimize \eqref{mcind:eq:sourceorder} on
these same fibres to obtain
\begin{equation}\tag{MC7}\label{mcind:eq:classorder}
 \boxed{\frac{h_*}{2}W\preceq G_0\preceq G\preceq\kappa G_0,
                    \qquad G_0\preceq UW.}
\end{equation}
$W$ is only an auxiliary matrix of the literal original $J$;
it is not substituted for either attained native metric.

The projection argument PD4--19 depends only on this bracket,
the fixed $K,A,x$, and the relative number $\kappa$. Here its
parameter is $\epsilon_{\rm PD}=2\theta$. This is not an
assertion that $2\theta$ equals PD1's different outer-truncation
expression $L_N^\sharp/(2J_{\rm trunc})$, or that its truncated
lower metric equals \eqref{mcind:eq:G}. We have supplied a new explicit
lower source minimum for the very same original $G,K,A,x$.

\section{The full complex mixed differences}

All inner products below are conjugate linear in the first
argument. Let $P_0,P$ be the $G_0,G$ orthogonal projections
onto the same $K$, and put
\[
 a=P_0x,\quad b=(I-P_0)x,\quad
 u=\mcindnorm{a}_{G_0},\quad v=\mcindnorm{b}_{G_0},\quad
 e=(P-P_0)x,\quad M_A=\mcindnorm{A}_{G_0}.
\]
For completeness the needed sharp projection estimate can be
proved entirely in this fixed class space. The operator
$E=G_0^{-1}G$ is self-adjoint in the $G_0$ inner product and
$I\preceq E\preceq\kappa I$ there. Decompose orthogonally as
$K\oplus K^{\perp_{G_0}}$ and write its operator blocks as
\[
 E=\begin{pmatrix}B&X\\X^{\dagger_0}&D\end{pmatrix},\qquad
 P=\begin{pmatrix}I&T\\0&0\end{pmatrix},\qquad T=B^{-1}X.
\]
The last formula is exactly the $G$-orthogonality condition for
the residual $(-Tb,b)$. Spectral calculus gives
$E^2\preceq(\kappa+1)E-\kappa I$. Its upper-left block and
congruence by $B^{-1}$ give
\[
 TT^{\dagger_0}\preceq(\kappa+1)B^{-1}-\kappa B^{-2}-I
 \preceq\frac{(\kappa-1)^2}{4\kappa}I.
\]
The scalar upper bound follows by maximizing
$(\kappa+1)/t-\kappa/t^2-1$ on $[1,\kappa]$; its critical
point is $2\kappa/(\kappa+1)$ and its maximum is the displayed
constant. Thus, without changing any original coordinate or mass,
\begin{equation}\tag{MC8}\label{mcind:eq:projection}
 \mcindnorm{e}_{G_0}\le sv,\qquad
 s=\frac{\kappa-1}{2\sqrt\kappa}
   =\frac{\theta}{\sqrt{1-2\theta}},\quad e\in K,
 \quad\mcindnorm{b-e}_{G_0}\le\sqrt{1+s^2}\,v.
\end{equation}
The last inequality uses $b\perp_{G_0}e$. If either summand
space is zero, the same statements follow directly with $e=0$.

Retain the four original complex pairings of PD3,
\begin{equation}\tag{MC9}\label{mcind:eq:pairings}
 \begin{array}{ll}
 z=\langle Px,A(I-P)x\rangle_G,&z_0=\langle a,Ab\rangle_{G_0},\\
 w=\langle(I-P)x,APx\rangle_G,&w_0=\langle b,Aa\rangle_{G_0}.
 \end{array}
\end{equation}
The exact differences, before any modulus or real part, are
\[
\begin{split}
 z-z_0={}&\langle e,Ab\rangle_{G_0}-\langle a,Ae\rangle_{G_0}
          -\langle e,Ae\rangle_{G_0}
          +\langle a+e,(E-I)A(b-e)\rangle_{G_0},\\
 w-w_0={}&\langle b,Ae\rangle_{G_0}-\langle e,Aa\rangle_{G_0}
          -\langle e,Ae\rangle_{G_0}
          +\langle b-e,(E-I)A(a+e)\rangle_{G_0}.
\end{split}
\]
Cauchy--Schwarz, \eqref{mcind:eq:projection}, and
$\mcindnorm{E-I}_{G_0}\le\kappa-1$ prove the unchanged full PD11
radius
\begin{equation}\tag{MC10}\label{mcind:eq:radius}
 |z-z_0|,|w-w_0|\le R(M_A),\quad
 R(M_A)=M_A\{sv^2+suv+s^2v^2
              +(\kappa-1)(u+sv)\sqrt{1+s^2}\,v\}.
\end{equation}
No term of $A=M_S$ is omitted or replaced by another action.

\section{Checking the proposed constant, and strengthening it}

Use exactly the proposed constants
\begin{equation}\tag{MC11}\label{mcind:eq:C}
 \overline M=2\sqrt{U/h_*}\,\mcindnorm{A}_W,\qquad
 C=\overline M U\mcindnorm{x}_W^2.
\end{equation}
The weaker $G_0\preceq2UW$ and lower bound in
\eqref{mcind:eq:classorder} already give $M_A\le\overline M$ by
applying the two norm comparisons to $Ay$ and $y$.
The stronger bound actually gives
$M_A\le\sqrt{2U/h_*}\mcindnorm{A}_W\le\overline M$; we keep
the original $\overline M$ unchanged. Orthogonality gives
$u^2+v^2=\mcindnorm{x}_{G_0}^2$. Consequently
\begin{equation}\tag{MC12}\label{mcind:eq:center}
 \mcindnorm{x}_{G_0}^2\le U\mcindnorm{x}_W^2,
 \qquad |z_0|,|w_0|\le M_Auv
                \le\tfrac12M_A\mcindnorm{x}_{G_0}^2\le C/2.
\end{equation}
In particular the proposed weaker $|z_0|,|w_0|\le C$ is valid.

For $0<\theta\le1/4$ we have
\[
 s\le\sqrt2\theta,\qquad s^2\le2\theta^2,\qquad
 \kappa-1\le4\theta,\qquad
 \sqrt{1+s^2}\le\frac3{2\sqrt2}.
\]
The coefficient of $uv$ in $R/M_A$ is at most
$\sqrt2\theta+4\theta\,3/(2\sqrt2)=4\sqrt2\theta$.
The coefficient of $v^2$ is at most
$\sqrt2\theta+2\theta^2+6\theta^2
\le(2+\sqrt2)\theta$. Hence
\begin{equation}\tag{MC13}\label{mcind:eq:Rsmall}
 R(M_A)\le M_A\theta\{4\sqrt2uv+(2+\sqrt2)v^2\}
          \le6M_A\theta(u^2+v^2)\le6C\theta.
\end{equation}
The middle inequality is proved by the exact positive difference
\[
 6(u^2+v^2)-4\sqrt2uv-(2+\sqrt2)v^2
 =6(u-\sqrt2v/3)^2+(8/3-\sqrt2)v^2\ge0.
\]
For comparison, the same coefficient bounds also prove every
step of the proposed original chain:
\[
 R(M_A)\le8M_A\theta v(u+v)
 \le12M_A\theta\mcindnorm{x}_{G_0}^2\le24C\theta.
\]
Indeed both coefficients are at most $8$, and
$uv+v^2\le\tfrac32(u^2+v^2)$. The last step uses only the
weaker $G_0\preceq2UW$, so that chain has no failed constant.

Put $\Delta z=z-z_0$, $\Delta w=w-w_0$. The exact complex
product, with every conjugation and remainder, is
\[
 \overline z w-\overline{z_0}w_0
  =\overline{\Delta z}\,w_0+\overline{z_0}\,\Delta w
                          +\overline{\Delta z}\,\Delta w.
\]
Using the proposed weaker bounds would give
$48C^2\theta+576C^2\theta^2\le192C^2\theta$, as requested.
Using \eqref{mcind:eq:center} and \eqref{mcind:eq:Rsmall} proves the stronger
same-constant result
\begin{equation}\tag{MC14}\label{mcind:eq:product}
 \boxed{|\overline z w-\overline{z_0}w_0|
       \le6C^2\theta+36C^2\theta^2\le15C^2\theta.}
\end{equation}
Only the last step uses $\theta^2\le\theta/4$. This encloses
the full complex product, hence its real and imaginary parts,
without assigning an unknown phase or sign.

\section{An explicit original moment tolerance and its scalar inputs}

For $C>0$ and any prescribed $t_*>0$, choose
\begin{equation}\tag{MC15}\label{mcind:eq:final}
 \boxed{\theta=\min\{1/4,t_*/(15C^2)\},\qquad
 0\le\varepsilon\le\theta h_*/d_{\rm src}.}
\end{equation}
Equations \eqref{mcind:eq:momentbudget}--\eqref{mcind:eq:product} then give
$|\overline z w-\overline{z_0}w_0|\le t_*$ with exact midpoint
arithmetic. The proposed denominator $192$ is also valid, but
needlessly smaller as a moment-accuracy allowance. If $C=0$,
positive definiteness of $W$ and positivity of $U/h_*$ imply
$x=0$ or $A=0$; all four pairings vanish. One may choose any
$\theta\le1/4$ to construct a positive midpoint metric, with
zero product error. If $v=0$, the actual sharper radius vanishes
as well: $x\in K$, both projections fix it, and all four pairings
vanish even when $C>0$.

The moment input in \eqref{mcind:eq:final} can be supplied by the
original onefold scalar integrals, not a proxy sequence. Write
$\eta_l=\int y^lw_h(y)dy$, with odd moments zero by evenness.
All source entries in \eqref{mcind:eq:H} use convolution moments
$\zeta_{2r}$ with $r\le N$, and
\[
 \zeta_L=\sum_{l_1+\cdots+l_k=L}
             \frac{L!}{l_1!\cdots l_k!}\prod_{j=1}^k\eta_{l_j},
 \qquad \int x^r d\nu_\epsilon=\zeta_{2(r+\epsilon)}.
\]
Absolute integrability follows from the original exponential
moment, so Fubini and the finite multinomial identity justify
these equalities, including all odd summands before they vanish.
In particular $\zeta_0=\mu_h^k$ is not replaced by one.
Let
\[
 V_N=\max_{0\le j\le N}(2j)!\alpha^{-2j}M_\alpha.
\]
The inequality $\cosh(\alpha y)\ge(\alpha y)^{2j}/(2j)!$
gives $0\le\eta_{2j}\le V_N$. Real certified midpoints for
these actual integrals, each with error at most
\begin{equation}\tag{MC16}\label{mcind:eq:onefold}
 \tau\le\min\left\{1,
         \frac{\varepsilon}{k^{2N+1}(V_N+1)^{k-1}}\right\},
\end{equation}
give every required convolution midpoint at error at most
$\varepsilon$. To prove this, telescope each $k$-fold product;
its error is at most $k\tau(V_N+1)^{k-1}$. The multinomial
coefficient sum at degree $L\le2N$ is $k^L\le k^{2N}$,
which proves \eqref{mcind:eq:onefold} after retaining the even terms.
This includes the native mass and the common parity correlation.

The integrals $\eta_{2j}$, $\vartheta_h$, $M_\alpha$ and the
original TW7/BT12 constants remain unevaluated original inputs.
Their compact-integral certification and original tail bounds
are specified in \cite{mcind:NP}; they are not produced by this finite
matrix argument. The matrices $J,A$ and the fixed class $x$ are
likewise their original data. If desired, the finite exact bound
\[
 \mcindnorm{A}_W\le\sqrt{\mcindtr(W^{-1}A^*WA)}
\]
can be used in \eqref{mcind:eq:C}; it follows by bounding the largest
eigenvalue of the positive matrix
$(W^{1/2}AW^{-1/2})^*(W^{1/2}AW^{-1/2})$ by its trace.
Using any certified upper bound here enlarges $C$ safely; it does
not change the actual action or select a different class.

The calculation controls absolute product error at the fixed
original degree. It does not assert a sign for its unevaluated
center, normalize a denominator, or establish a growing-degree
phase limit. Any later metric-dependent denominator must keep its
own original enclosure. Numerical evaluation errors in $z_0,w_0$
must also be added: if their centers have errors $\delta_z,
\delta_w$, the extra product radius is at most
$(C/2)(\delta_z+\delta_w)+\delta_z\delta_w$. This follows
from the same three-term complex expansion and is not silently
absorbed into exact midpoint arithmetic.

\begin{thebibliography}{9}
\bibitem{mcind:PD} Split-Zero programme,
\emph{Sharp original-metric projection control and disks for complex
mixed pairings}, 20 September 2026, complete source
\path{projection_intake_20260920/PROJECTION_PHASE_RETURN.tex},
PD1--23. Read completely for this independent check.
\bibitem{mcind:NP} Split-Zero programme,
\emph{Explicit native moment precision, inverse margins, and the
original signed receiver}, 20 September 2026, complete source
\path{independent_native_moment_precision/NATIVE_MOMENT_PRECISION.tex},
especially equations (5), (13), (39)--(43), (51), (53)--(56).
This is the preceding independent native calculation, not an
assumed positivity oracle. Its exact rational test instances are
not the original arithmetic measure.
\bibitem{mcind:CAI} Split-Zero programme,
\emph{The complete original action relative to its arithmetic total},
complete provider \path{CAI_COMPLETE_PREREQUISITE.tex}, CAI16--17,
retaining the original TW7 and BT12 constants and full native
Gamma envelope. The envelope is used for bounds only; the
measure remains $w_h^{*k}$ throughout.
\end{thebibliography}

\endgroup

% END COMPLETE PRECISION TRANSFER AND PHASE PROOFS
